\maketitle
\noteprovenance

\begin{abstract}
The three-prime reciprocal running-LCM kernel has nonsingular minors of
every order, uniformly in its third-coordinate layer.  A binary floor carry
explains its non-separability, exact sampled ranks and sharp uniform
approximation obstruction.  This is compatible with finite-rank
approximation of the original kernel in the summation norm.
For the repeated $\{2,3,5\}$ series we derive the literal integer tail
recurrence, two quadratic bounds, exact denominator clearing and a
residue-escape equivalence for dominating caps of size $o(8^a)$.
We sharpen the fixed-start residue description and isolate the moving
boundary in a value-preserving weighted difference; a literal counterexample
shows that absence of floor crossings alone does not force cancellation.
No cofinal escape or irrationality theorem is proved for this series.
The two-prime comparison reproduces the earlier reduction attributed to
Steve Fan, using the classical Hecke--Mahler value theorem.
Ordinary proofs, external analytic inputs, supplied formal receipts and
finite computations are kept separate.
\end{abstract}

\paragraph{How the short paper uses this record.}
The short paper centres on the binary carry and its rank obstruction.
Section~\ref{long269:sec:rank} gives the rank and approximation arguments;
Section~\ref{long269:sec:two-prime} gives the two-prime comparison.
Sections~\ref{long269:sec:blocks}--\ref{long269:sec:escape} recover the
literal tails, denominator clearing and window criterion, while
Section~\ref{long269:sec:evidence} retains the finite certificates.
These are two distinct mathematical tasks: the rank argument excludes a
separated representation, and the tail argument identifies the arithmetic
condition still needed for irrationality.  Neither substitutes for the other.

\section{The problem, and what is settled}\label{long269:sec:problem}

Let $P$ be a finite set of primes with $|P|\ge2$, and let
$a_1<a_2<\cdots$ enumerate the positive integers all of whose prime factors lie
in $P$.  Erd\H{o}s Problem~\#269 asks whether
\[
 \sum_{n\ge1}\frac{1}{[a_1,\ldots,a_n]}
\]
is irrational, where $[a_1,\ldots,a_n]$ is the least common
multiple~\cite[p.~65]{erdosgraham1980}\cite[p.~106]{erdos1988}.  Bloom's
catalogue is the source of the problem number and the status recorded in
the cited snapshot~\cite{erdosproblems}.  This revision does not make a
new claim about the live catalogue status.
Theorem~\ref{long269:res:lead-two-prime} settles every instance with $|P|=2$,
at the level of transcendence, and Fan posted that argument
first~\cite{fan2026comment}; this record leaves the repeated finite cases with $|P|\ge3$ unresolved.

The restriction $|P|\ge2$ is necessary.  For $P=\{p\}$ the enumeration is
$a_n=p^{\,n-1}$, so $[a_1,\ldots,a_n]=p^{\,n-1}$ and the sum is $p/(p-1)$.  For
infinite $P$ the sum is always irrational, which Erd\H{o}s calls a simple
exercise~\cite[p.~106]{erdos1988}.

Write $\mathcal R_P$ for the sum above and $\mathcal D_P$ for the
de-duplicated sum, in which each distinct value of the running least common
multiple contributes its reciprocal once.  The two differ because the running
value is constant along stretches of the enumeration.  In a letter written on
1~January 1973 Erd\H{o}s recorded that he could prove irrationality of the
de-duplicated sum~\cite[p.~335]{erdos1974letter}.  He states that assertion for
given primes $p_1,\ldots,p_r$ and supplies no proof; for a single prime the
de-duplicated sum is again $p/(p-1)$, so the assertion concerns sets of at
least two primes.  The open question addressed by the catalogue is
$\mathcal R_P$; its two-prime case is Theorem~\ref{long269:res:lead-two-prime}.

\begin{theorem}[two-prime transcendence]\label{long269:res:lead-two-prime}
Let $p$ and $q$ be distinct primes.  Then $\mathcal R_{\{p,q\}}$ and
$\mathcal D_{\{p,q\}}$ are transcendental.
\end{theorem}

Section~\ref{long269:sec:two-prime} proves this.  Both values are nonconstant polynomials over $\mathbb Q$
in one Hecke--Mahler boundary sum $A$,
and transcendence of $A$ is the theorem of Loxton and van der Poorten
\cite[Theorem~8, p.~40]{loxtonvdp1977} in the modern form of Bugeaud and
Laurent \cite[Theorem~1.1]{bugeaudlaurent2023}.  The evidence class
is an ordinary proof in this paper over a cited external theorem, and no Lean
declaration formalises it.

\paragraph{Attribution.}
The supplied forum record credits Steve Fan's post of 26 June 2026 with
the running-LCM identity for finite prime sets, the two-channel
factorisation, the Hecke--Mahler reduction and its transcendence conclusion
\cite{fan2026comment}.  The unrestricted analogue is classical: iterating
the prime-exponent maximum rule gives
$\operatorname{lcm}(1,\ldots,N)=\prod_{t\le N}t^{\lfloor\log_t N\rfloor}$
over the primes $t\le N$, equivalently
$\log\operatorname{lcm}(1,\ldots,N)=\psi(N)$; see Apostol~\cite{apostol1976}
and Montgomery and Vaughan~\cite{montgomeryvaughan2007}.  We reproduce that
argument to fix normalisations and distinguish repeated from distinct-height
sums, not as a new application of the value theorem.  The live thread was inaccessible during this
revision; priority is retained from the supplied record rather than
re-presented as independently verified correspondence.  The scalar
transcendence theorem itself is credited separately to its published authors.

\paragraph{What the third prime does.}
Section~\ref{long269:sec:rank} isolates the obstruction exactly.  After separating row
and column factors, the two-prime kernel is an outer product and the
three-prime kernel retains one binary floor carry.  That carry is enough to
produce nonsingular minors of every order, uniformly in the third coordinate,
and to make the normalised carry matrix inapproximable by finite separated sums
below half its own jump.  The minors use only row and column rescaling, the
density of the rotations by $\log_rp$ and $\log_rq$ taken separately, and a
staircase determinant, and one choice of indices serves every layer of the
literal kernel.  Sections~\ref{long269:sec:blocks} to~\ref{long269:sec:escape} take a
second route on the literal $\{2,3,5\}$ series: dyadic shells give an
integer-coefficient tail recurrence, a strict endpoint clears the smooth part
of a hypothetical denominator, and the surviving carry is trapped between a
quadratic bound and a residue condition.  Section~\ref{long269:sec:evidence} reports
the finite evidence and its ceiling, and Section~\ref{long269:sec:open} states what
remains.

Throughout, $p,q,r$ are pairwise distinct primes.  Call $n$ \emph{smooth} when
$n=p^{i}q^{j}r^{k}$ for some $i,j,k\ge0$; this is the
smooth lattice value.  For $x\ge1$ write
\[
 \Lprefix(x)=\lcm\{\,n\le x:\ n\ \text{smooth}\,\},
 \qquad
 \hgt(x)=p^{\lfloor\log_p x\rfloor}\,q^{\lfloor\log_q x\rfloor}\,
 r^{\lfloor\log_r x\rfloor},
\]
the
running least common multiple and the
pure-power height, where $\lfloor\log_b x\rfloor$ is the largest $e$ with $b^{e}\le x$.
Since $a_1,\ldots,a_n$ are exactly the smooth numbers up to $a_n$, we have
$[a_1,\ldots,a_n]=\Lprefix(a_n)$.  The reciprocal of the height at a smooth
point is the
lattice kernel
\[
 \Kernel(i,j,k)=\frac{1}{\hgt(p^{i}q^{j}r^{k})} .
\]
Here ``smooth'' always means supported on the fixed prime set, and it is not
the varying-bound notion counted by $\Psi(x,y)$ in the Dickman--Hildebrand
theory \cite{hildebrand1986}; no smooth-number density asymptotic is
used below.  The appropriate fixed-support context is Tijdeman--Meijer
\cite[Sections~1--4]{tijdemanmeijer1974}.  The modern two-prime treatment
of Languasco, Luca, Moree and Togb\'e \cite{languasco2025} also makes the
lattice-triangle geometry explicit; its gap estimates are not inputs to
our shell bound.  For $b\in\{p,q,r\}$ the set of positive powers of $b$ is the
\emph{$b$-channel}.

The statement of Problem~\#269 has been formalised before, as a conjecture with
an unfilled proof, in the \emph{Formal Conjectures}
collection~\cite{formalconjectures269}.  That is a formal statement of the
question up to a rational normalisation, since its Nat-indexed series includes
the empty-prefix term, and its rational, irrational and infinite-prime
assertions all end in \texttt{sorry}.  Kova\v{c} and Tao~\cite{kovactao2024}
treat several irrationality problems of Erd\H{o}s for series of unit fractions
by elementary means; nothing from that work is used here.

\medskip
\noindent\textbf{Keywords.} irrationality; transcendence; least common
multiple; smooth numbers; separated rank; Lean~4.
\qquad
\textbf{MSC 2020.} 11J72 (primary); 11A05, 11N25, 68V20 (secondary).

\subsection*{Notation and proof dependencies}
There are two different sums: $\mathcal R_P$ counts every supported integer,
whereas $\mathcal D_P$ counts each height once.  From the dyadic section
onwards, $S=\mathcal R_{\{2,3,5\}}$.  Write
$P_a=\hgt(2^a)$, $h_a=P_a/2$, $s_a$ for the shell mass and
$X_a=h_a\sum_{j\ge a}s_j$.  Thus $P_a$ is a boundary height, not a set of
primes.  The long record's $U_a$ is the short note's $T_a$.
The cumulative jump count is
$n_a=a+\lfloor a\log_3 2\rfloor+\lfloor a\log_5 2\rfloor$;
the strict index is $j_a=n_a-1$ for $a\ge1$.
The long cap is $K(B,a)=\lfloor BQ(n_a)\rfloor$, while the short note
uses $K_0(B,a)=90B(a+1)^2$.  Neither is a residue or a hypothetical
integer carry.  A window has start $\ell$, length $h$, product $W_{\ell,h}$
and forcing $F_{\ell,h}$; zero modulo $W$ is represented by $W$, not zero.

The rank proof uses only diagonal rescaling and separate one-dimensional
density.  The arithmetic proof instead uses the actual shell coefficients:
summability $\to$ recurrence $\to$ denominator clearing $\to$ the residue
criterion.  The weighted-shift discussion is a third route, with explicit
unproved cancellation hypotheses.  No arrow from kernel rank to scalar
irrationality is asserted.


\section{The finite geometry of the running value}\label{long269:sec:lcm}

The smooth numbers up to $x$ are indexed by the exponent triples $(i,j,k)$ with
$i\le\lfloor\log_p x\rfloor$, $j\le\lfloor\log_q x\rfloor$,
$k\le\lfloor\log_r x\rfloor$ and $p^{i}q^{j}r^{k}\le x$, the
smooth prefix index set.  Both conditions matter: the coordinate box is strictly
larger than the prefix, since a product of three large pure powers can exceed
$x$ while each factor does not.

\begin{theorem}[the running least common multiple]\label{long269:res:lcm}
Let $p,q,r$ be pairwise distinct primes and $x\ge1$.  Then
$\Lprefix(x)=\hgt(x)$.
\end{theorem}

\begin{proof}
Every smooth $n\le x$ has exponents bounded by the corresponding integer
logarithms, so $n\mid\hgt(x)$ and hence $\Lprefix(x)\mid\hgt(x)$.  Conversely
the three pure powers $p^{\lfloor\log_p x\rfloor}$,
$q^{\lfloor\log_q x\rfloor}$ and $r^{\lfloor\log_r x\rfloor}$ are themselves
smooth numbers not exceeding $x$, so each divides $\Lprefix(x)$, and distinct
primes have coprime powers, so their product divides $\Lprefix(x)$ as well.
\end{proof}



At $(p,q,r)=(2,3,5)$ the first ten values are
\[
\begin{array}{c|cccccccccc}
x&1&2&3&4&5&6&7&8&9&10\\ \hline
\Lprefix(x)&1&2&6&12&60&60&60&120&360&360
\end{array}
\]
So $\Lprefix(6)=4\cdot3\cdot5=60$, which exceeds $6$: the running value at a
smooth cutoff already contains powers of the other two primes that the cutoff
itself does not.  That small fact is the mechanism behind
Section~\ref{long269:sec:rank}.  Also $\hgt(x)\le x^{3}$, since each of the three
factors is a power of its base not exceeding $x$; this is the
cubic majorant, and the exponent is the number of generating primes.

Say that $x$ and $y$ lie in the same \emph{logarithmic cell} when
$\lfloor\log_b x\rfloor=\lfloor\log_b y\rfloor$ for each of $b=p,q,r$, the
cell relation.  By Theorem~\ref{long269:res:lcm} the running value depends on $x$ only
through the three integer logarithms, so it is constant on cells and moves only
where one logarithm moves.

\begin{proposition}[constancy and jump ratios]\label{long269:res:cell}
If $x,y\ge1$ lie in the same logarithmic cell then $\Lprefix(x)=\Lprefix(y)$,
and the same holds for the kernel at two smooth points of one cell.  If
$\lfloor\log_p y\rfloor=\lfloor\log_p x\rfloor+1$ while the other two
logarithms agree, then $\Lprefix(y)=p\,\Lprefix(x)$, and similarly with $q$ or
$r$ in place of $p$.
\end{proposition}

\begin{proof}
Both parts are immediate from Theorem~\ref{long269:res:lcm}: the height depends on $x$
only through the three integer logarithms, and advancing one of them multiplies
exactly one factor by its base.
\end{proof}

\label{long269:res:jump}

\begin{proposition}[jump count]\label{long269:res:count}
Let $n\ge0$.  The set of the first $n$ positive powers of $p$, of $q$ and of
$r$ has exactly $3n$ elements, and adjoining the common origin $1$ gives
exactly $3n+1$.
\end{proposition}

\begin{proof}
Within one channel the powers $b,b^{2},\ldots,b^{n}$ are distinct because
$b\ge2$.  Across two channels a common value would be a positive power of two
distinct primes, which unique factorisation forbids.  Finally $1$ is not a
positive power of any prime.
\end{proof}



Two finite statements are recorded here because they are the exact shape of the
counting used later.  Write $\mathcal B(h_p,h_q,h_r)$ for the box of exponent
triples with $i\le h_p$, $j\le h_q$, $k\le h_r$, the
exponent box, and $F(H)$ for the set of its points whose height equals $H$, the
height fibre.

\begin{proposition}[finite normal form]\label{long269:res:fibre}
For every box $\mathcal B$,
$\sum_{(i,j,k)\in\mathcal B}\Kernel(i,j,k)=\sum_{H}\#F(H)/H$, the outer sum
ranging over the heights attained on $\mathcal B$.
\end{proposition}

\begin{proof}
Partition $\mathcal B$ into the fibres of the height map.  On $F(H)$ every
summand is $1/H$, so the fibre contributes $\#F(H)/H$.
\end{proof}



Now fix an interval $[\lo,\hi)$ and write $\mathcal S$ for the exponent triples
of $\mathcal B(h_p,h_q,h_r)$ whose smooth value lies in it, the
smooth exponent shell.

\begin{lemma}[uniqueness in a short interval]\label{long269:res:short}
Let $b\ge1$ and $\hi\le b\,\lo$.  If $b^{a}w$ and $b^{a'}w$ both lie in
$[\lo,\hi)$ then $a=a'$.
\end{lemma}

\begin{proof}
If $a<a'$ then $\hi\le b\,\lo\le b^{a+1}w\le b^{a'}w<\hi$, which is impossible;
the case $a>a'$ is symmetric.
\end{proof}

\begin{proposition}[projection and the quadratic shell bound]\label{long269:res:drop}
\label{long269:res:shell}
If $\hi\le r\,\lo$ then $\#\mathcal S\le(h_p+1)(h_q+1)$, and if
$\hi\le p\,\lo$ then $\#\mathcal S\le(h_q+1)(h_r+1)$.  If moreover
$\hi\le r\,\lo$ and $h_p\le h_q\le h_r$ with $h_p+h_q+h_r=j$, then
$9\,\#\mathcal S\le(j+3)^{2}$.
\end{proposition}

\begin{proof}
Suppose $\hi\le r\,\lo$.  If two triples of $\mathcal S$ agree in their first
two coordinates, Lemma~\ref{long269:res:short} with $b=r$ and $w=p^{i}q^{j}$ forces
their third coordinates to agree, so the projection forgetting the third
coordinate is injective on $\mathcal S$ and its image lies in a rectangle with
$(h_p+1)(h_q+1)$ points.  The other case is the same with the first coordinate
projected away.  Under the sorting hypothesis the two surviving coordinates are
the two smallest, so it suffices that $a\le b\le c$ with $a+b+c=j$ gives
$9(a+1)(b+1)\le(j+3)^{2}$.  From $a\le b\le c$ we get $a+2b\le j$, so it is
enough that $9(a+1)(b+1)\le(a+2b+3)^{2}$; writing $b=a+d$ with $d\ge0$, the
difference of the two sides is $d(3a+4d+3)\ge0$.
\end{proof}



\section{One Hecke--Mahler value controls both two-prime sums}
\label{long269:sec:two-prime}

Temporarily let $P=\{p,q\}$ with $p<q$, and write
$L_{p,q}(t)=p^{\lfloor\log_p t\rfloor}q^{\lfloor\log_q t\rfloor}$, which is the
running least common multiple of the $\{p,q\}$-smooth numbers up to $t$ by the
argument of Theorem~\ref{long269:res:lcm} with one coordinate omitted.  The
de-duplicated sum retains the initial value $1$ and one reciprocal for every
later distinct running value, so
\[
 \mathcal D_{\{p,q\}}
 =1+\sum_{t\in\{p,p^2,\ldots\}\cup\{q,q^2,\ldots\}}\frac1{L_{p,q}(t)},
 \qquad
 \mathcal R_{\{p,q\}}=\sum_{i,j\ge0}\frac1{L_{p,q}(p^iq^j)} .
\]

\begin{proof}[Proof of Theorem~\ref{long269:res:lead-two-prime}]
Set
\[
 \theta=\frac{\log p}{\log q},\qquad x=\frac1p,\qquad y=\frac1q,\qquad
 m_n=\lfloor n\theta\rfloor,\qquad \delta_n=m_{n+1}-m_n .
\]
Here $0<\theta<1$, and $\theta$ is irrational, since a rational value would
give $p^{b}=q^{a}$ for positive integers $a,b$.  Consequently
$\delta_n\in\{0,1\}$.  Put
\[
 A=\sum_{n\ge0}x^ny^{m_n},
 \qquad
 B_\ast=\sum_{n\ge0}\delta_nx^ny^{m_n+1} .
\]
The initial value and the $p$-channel contribute $A$, since
$L_{p,q}(p^n)=p^nq^{m_n}$.  A power of $q$ lies strictly between $p^n$ and
$p^{n+1}$ exactly when $\delta_n=1$, it is then $q^{m_n+1}$, and its post-jump
reciprocal is $x^ny^{m_n+1}$; so the $q$-channel contributes $B_\ast$ and
$\mathcal D_{\{p,q\}}=A+B_\ast$.  All these series converge absolutely.

Since $y^{m_{n+1}}-y^{m_n}=\delta_ny^{m_n}(y-1)$, an index shift gives
$A-1-xA=x(y-1)B_\ast/y$, so $B_\ast=\bigl((p-1)A-p\bigr)/(1-q)$ and
\begin{equation}\label{long269:eq:two-prime-affine}
 \mathcal D_{\{p,q\}}=\frac{(q-p)A+p}{q-1} .
\end{equation}
At a smooth point, $\log_p(p^iq^j)=i+j/\theta$ and
$\log_q(p^iq^j)=j+i\theta$, so
$L_{p,q}(p^iq^j)=p^{\,i+\lfloor j/\theta\rfloor}q^{\,j+m_i}$ and absolute
convergence permits the factorisation
\[
 \mathcal R_{\{p,q\}}=AC,\qquad C=\sum_{j\ge0}y^jx^{\lfloor j/\theta\rfloor}.
\]
For $j\ge1$ put $n=\lfloor j/\theta\rfloor$.  Then $n>j/\theta-1$ gives
$n\theta>j-\theta>j-1$, and $n+1>j/\theta$ gives $(n+1)\theta>j$, while
$n\theta\le j$ and $(n+1)\theta<j+\theta<j+1$; hence $m_n=j-1$, $m_{n+1}=j$ and
$\delta_n=1$.  Conversely every $n$ with $\delta_n=1$ arises from the unique
$j=m_{n+1}$.  Therefore $C=1+B_\ast$ and
\begin{equation}\label{long269:eq:two-prime-quadratic}
 \mathcal R_{\{p,q\}}
 =\frac{(p+q-1)A-(p-1)A^{2}}{q-1} .
\end{equation}

It remains to prove that $A$ is transcendental.  For the Hecke--Mahler series
\[
 F_\theta(x,y)=\sum_{n\ge1}\sum_{k=1}^{\lfloor n\theta\rfloor}x^ny^k
\]
a finite geometric sum gives
$\bigl((1-y)/y\bigr)F_\theta(x,y)=x/(1-x)-(A-1)$, that is
\begin{equation}\label{long269:eq:hecke-mahler-boundary}
 A=\frac1{1-x}-\frac{1-y}{y}F_\theta(x,y).
\end{equation}
Bugeaud and Laurent's Theorem~1.1 states, in particular, that
$F_\theta(\beta,\alpha)$ is transcendental when $\theta\in(0,1)$ is irrational,
$\alpha$ and $\beta$ are nonzero algebraic numbers, $|\beta|<1$ and
$|\beta\alpha^\theta|<1$ \cite[Theorem~1.1]{bugeaudlaurent2023}; the
$\rho=0$ case used here goes back to Loxton and van der Poorten
\cite[Theorem~8, p.~40]{loxtonvdp1977}.  Take $(\beta,\alpha)=(x,y)$: then
$|\beta|=1/p<1$ and
\[
 |xy^{\theta}|=\frac1p\left(\frac1q\right)^{\log p/\log q}=\frac1{p^{2}}<1 .
\]
So $F_\theta(x,y)$ is transcendental, and \eqref{long269:eq:hecke-mahler-boundary}
makes $A$ transcendental.  The coefficient of $A$
in~\eqref{long269:eq:two-prime-affine} is $(q-p)/(q-1)\ne0$ and the coefficient of
$A^{2}$ in~\eqref{long269:eq:two-prime-quadratic} is $-(p-1)/(q-1)\ne0$, both rational.
If either value were algebraic, its identity would exhibit $A$ as a root of a
nonzero polynomial over the algebraic numbers.
\end{proof}

\noindent The two identities are separately labelled
\label{long269:res:two-prime-transcendence}\label{long269:res:two-prime-repeated-transcendence}
because they carry different weight.  A product of two transcendental numbers
need not be transcendental, so the factorisation $\mathcal R=AC$ alone proves
nothing; the quadratic identity~\eqref{long269:eq:two-prime-quadratic} in the single
value $A$ is what settles the repeated sum.  A third prime replaces the single
Beatty boundary by a genuinely two-dimensional ordering problem, and the next
section makes that failure exact.

\section{The rank phase transition}\label{long269:sec:rank}

One might hope to write the kernel as $f(i)g(j)h(k)$ and reduce the problem to
one-dimensional criteria.  At two generators that hope is exactly correct, and
at three it fails at every finite order.

\begin{proposition}[two generators separate]\label{long269:res:two-prime-rank}
For real $p,q>1$, with $L_{p,q}(t)=p^{\lfloor\log_p t\rfloor}q^{\lfloor\log_q t\rfloor}$
as above, and all integers $i,j\ge0$, the two-prime kernel
$\Kernel_2(i,j)=1/L_{p,q}(p^iq^j)$ is the outer product
\[
 \Kernel_2(i,j)
 =\bigl(p^{i}q^{\lfloor\log_q p^{i}\rfloor}\bigr)^{-1}
  \bigl(p^{\lfloor\log_p q^{j}\rfloor}q^{j}\bigr)^{-1},
\]
so every two-by-two minor of $\Kernel_2$ vanishes.
\end{proposition}

\begin{proof}
Since $i$ and $j$ are integers,
$\lfloor\log_p(p^iq^j)\rfloor=i+\lfloor\log_p q^{j}\rfloor$ and
$\lfloor\log_q(p^iq^j)\rfloor=j+\lfloor\log_q p^{i}\rfloor$.  Hence
$L_{p,q}(p^iq^j)$ is the product of $p^{i}q^{\lfloor\log_q p^{i}\rfloor}$,
which depends on $i$ alone, and $p^{\lfloor\log_p q^{j}\rfloor}q^{j}$, which
depends on $j$ alone.  A matrix whose entries are a product of a row function
and a column function has vanishing two-by-two minors.  For distinct primes
$p,q$ the product $L_{p,q}$ is the running least common multiple by the
argument of Theorem~\ref{long269:res:lcm}.
\end{proof}



At three generators the smallest rectangle already fails to factor.  A
factorisation $f(i)g(j)h(k)$ would force
$\Kernel(0,0,0)\Kernel(1,1,0)=\Kernel(1,0,0)\Kernel(0,1,0)$.

\begin{proposition}[non-separability at $\{2,3,5\}$]\label{long269:res:rank}
With $(p,q,r)=(2,3,5)$,
\[
 \det\begin{pmatrix}
 \Kernel(0,0,0)&\Kernel(0,1,0)\\
 \Kernel(1,0,0)&\Kernel(1,1,0)
 \end{pmatrix}
 =\det\begin{pmatrix}1&1/6\\1/2&1/60\end{pmatrix}
 =-\frac1{15}\ne0 .
\]
\end{proposition}

\begin{proof}
The four values are computed from $\hgt(1)=1$, $\hgt(2)=2$,
$\hgt(3)=2\cdot3=6$ and $\hgt(6)=4\cdot3\cdot5=60$, so the determinant is
$1/60-1/12=-1/15$.
\end{proof}



\begin{theorem}[arbitrary-order non-separability]\label{long269:res:infinite-rank}
\label{long269:res:lead-infinite-rank}
Let $p,q,r$ be primes with $p\ne q$, $p\ne r$ and $q\ne r$.  For every
$n\ge1$ there are injective maps $I,J:\{0,\ldots,n-1\}\to\N$ such that, for
every $k\ge0$,
\[
 \det\bigl(\Kernel(I(a),J(b),k)\bigr)_{0\le a,b<n}\ne0 .
\]
Consequently, for no finite $d$ do there exist rational-valued functions
$f_\ell:\N\to\Q$ and $G_\ell:\N^{2}\to\Q$, $0\le\ell<d$, satisfying
$\Kernel(i,j,k)=\sum_{\ell<d}f_\ell(i)G_\ell(j,k)$ for all $i,j,k$.
\end{theorem}

\begin{proof}
Put $\alpha=\log_r p$, $\beta=\log_r q$, $x_i=\{i\alpha\}$ and
$y_j=\{j\beta\}$.  The three height exponents are
\[
 \begin{split}
 v_p(\hgt(p^iq^jr^k))&=i+\lfloor j\log_pq+k\log_pr\rfloor,\\
 v_q(\hgt(p^iq^jr^k))&=j+\lfloor i\log_qp+k\log_qr\rfloor,\\
 v_r(\hgt(p^iq^jr^k))&=k+\lfloor i\alpha\rfloor+\lfloor j\beta\rfloor+
   \lfloor x_i+y_j\rfloor .
 \end{split}
\]
Every term except the last floor depends on $i$ and $k$ alone or on $j$ and $k$
alone, so with positive rational $R_i(k)$ and $C_j(k)$,
\begin{equation}\label{long269:eq:carry-factorisation}
 \Kernel(i,j,k)=R_i(k)\,C_j(k)\,t^{\lfloor x_i+y_j\rfloor},
 \qquad t=r^{-1} .
\end{equation}
The remaining matrix is independent of $k$, and $x_i+y_j\in[0,2)$, so its
entries are $1$ and $t$.

Each of $\alpha$ and $\beta$ is irrational, since a rational value would give
an equality of positive powers of distinct primes, so each fractional-part
orbit is dense in $(0,1)$ and each is injective.  No simultaneous density of
the pair of rotations is required.  First choose indices with
$0<x_{I(0)}<\cdots<x_{I(n-1)}<1$.  For the columns write
$s_b=1-y_{J(b)}$, and use density of the second orbit to choose
\[
 s_0\in(0,x_{I(0)}),\qquad
 s_b\in(x_{I(b-1)},x_{I(b)})\quad(1\le b<n).
\]
The intervals are disjoint, so the $J(b)$ are distinct.  Their strict
endpoints ensure that no selected entry lies on a threshold, and
$x_{I(a)}+y_{J(b)}\ge1$ holds exactly when $b\le a$.  Dividing row $a$ by
$R_{I(a)}(k)$ and column $b$ by $C_{J(b)}(k)$ therefore leaves
\[
 C_n(t)=\begin{pmatrix}
 t&1&1&\cdots&1\\
 t&t&1&\cdots&1\\
 \vdots&\vdots&\ddots&\ddots&\vdots\\
 t&t&\cdots&t&1\\
 t&t&\cdots&t&t
 \end{pmatrix},
 \qquad
 \det C_n(t)=t(t-1)^{n-1}\ne0 ,
\]
the determinant following by subtracting from each row its predecessor, working
upwards from the last.  Before the division, the determinant equals
\[
 \det C_n(t)\prod_{a<n}R_{I(a)}(k)\prod_{b<n}C_{J(b)}(k),
\]
which is nonzero for every $k$; this also explains why the same indices work
in every layer.

For the last assertion, suppose that a representation with $d$ summands
exists and fix $k$.  On the selected rows $I(0),\ldots,I(d)$ and columns
$J(0),\ldots,J(d)$ the resulting matrix factors as
\[
 (f_\ell(I(a)))_{0\le a\le d,\,\ell<d}
 (G_\ell(J(b),k))_{\ell<d,\,0\le b\le d}.
\]
It has rank at most $d$ and hence zero determinant, contradicting the
nonzero minor of order $d+1$.
\end{proof}



Index selection is essential.  The leading minors of the
$\{2,3,5\}$ kernel are not a witness: row three is $1/120$ times row zero for
$j\le3$
(checked),
and the proportionality fails at $j=4$
(checked).
A proof that read off leading minors alone would be false.

The same threshold description determines every finite sampled rank, rather
than only producing one nonsingular minor.

\begin{proposition}[finite sampled cut rank]\label{long269:res:finite-cut-rank}
Let $m\ge1$ and let $c$ lie in a field with $c\ne0,1$.  For $0\le h\le m$,
let $v_h$ be the length-$m$ column whose first $h$ entries are $1$ and whose
remaining entries are $c$.  If the distinct columns of a matrix are the
$v_h$ with $h$ in a nonempty set $E\subseteq\{0,\ldots,m\}$, then its rank is
\[
 |E|-\mathbf 1_{\{0,m\}\subseteq E}.
\]
\end{proposition}

\begin{proof}
Write $E=\{h_1<\cdots<h_s\}$.  The $s-1$ consecutive differences are
\[
 v_{h_{a+1}}-v_{h_a}
   =(1-c)\mathbf 1_{\{h_a,\ldots,h_{a+1}-1\}},
\]
so they are linearly independent because their nonempty supports are
disjoint.  If $h_1>0$ or $h_s<m$, those supports miss a coordinate on which
$v_{h_1}$ is nonzero, and adjoining $v_{h_1}$ gives rank $s$.  If
$h_1=0$ and $h_s=m$, the differences sum to $(1-c)\mathbf 1$ while
$v_0=c\mathbf 1$, so $v_0$ is already in their span and the rank is $s-1$.
\end{proof}

To apply the proposition, sort any chosen row phases $x_i$.  Each normalised
column is a threshold column $v_h$, where $h$ is the number of sampled phases
strictly below $1-y_j$.  Repeated threshold positions give proportional
columns before column normalisation, and the nonzero row and column factors in
\eqref{long269:eq:carry-factorisation} do not change rank.  Thus the displayed
formula is the exact rank of every finite sample.

\paragraph{Attribution and reach.}
Fan recorded the qualitative obstruction for $|P|\ge3$ in his comment of
26 June 2026~\cite{fan2026comment}, observing that the two-prime argument does
not seem to generalise immediately.  The quantitative form above, its
uniformity in the third coordinate, and its formalisation are proved here.
Theorem~\ref{long269:res:infinite-rank} excludes exact separations of the displayed
finite-sum form; it is not an independence, irrationality or transcendence
statement, and it settles no instance of Erd\H{o}s~\#269.

\subsection{A sharp uniform obstruction for the normalised carry matrix}

A fixed nonzero determinant remains nonzero under sufficiently small
perturbations.  What the preceding argument does not provide is one error
threshold that works uniformly for the entire infinite matrix and all
orders.  The next theorem
replaces exactness by a uniform distance, and the constant it produces is
sharp.  It concerns the normalised carry matrix
\begin{equation}\label{long269:eq:carry-matrix}
 C(i,j)=t^{\lfloor x_i+y_j\rfloor},
 \qquad
 x_i=\{i\log_r p\},\quad y_j=\{j\log_r q\},\quad t=r^{-1},
\end{equation}
which is the factor left in~\eqref{long269:eq:carry-factorisation} after the row and
column factors are divided out.  Say that a real matrix $A$ on $\N\times\N$ has
\emph{finite separated rank} when all of its columns lie in one
finite-dimensional space of real sequences, equivalently when
$A(i,j)=\sum_{\ell<d}f_\ell(i)g_\ell(j)$ for some finite $d$ and some
sequences $f_\ell,g_\ell$, with no continuity or boundedness assumed.

\begin{theorem}[sharp uniform separated approximation]\label{long269:res:uniform-rank}
Let $p,q,r$ be pairwise distinct primes and let $C$ be as
in~\eqref{long269:eq:carry-matrix}.  Then
\[
 \inf_{A}\ \sup_{i,j\ge0}\ |C(i,j)-A(i,j)|=\frac{1-t}{2}=\frac{r-1}{2r},
\]
the infimum being over all matrices $A$ of finite separated rank, and it is
attained by the constant matrix of value $(1+t)/2$.
\end{theorem}

\begin{proof}
Every entry of $C$ lies in $\{1,t\}$, and $C(i,j)=t$ exactly when
$x_i+y_j\ge1$.  Fix $j\ne k$.  The numbers $y_j$ are pairwise distinct, since
$\log_r q$ is irrational, so we may assume $y_j<y_k$, and then
$0\le1-y_k<1-y_j\le1$.  Density of the orbit $(x_i)$ in $(0,1)$ supplies an
index $i$ with $1-y_k<x_i<1-y_j$, and at that row the two columns carry the
entries $t$ and $1$.  Hence any two distinct columns of $C$ are at sup-distance
exactly $1-t$.

Let $A$ have finite separated rank and put
$\varepsilon=\sup_{i,j}|C(i,j)-A(i,j)|$.  Suppose $\varepsilon<(1-t)/2$.  Each
column $A_j$ then satisfies $\|A_j\|_\infty\le1+\varepsilon$, so all columns of
$A$ lie in $W=V\cap\ell^{\infty}$, where $V$ is the finite-dimensional space
containing the columns of $A$; the space $W$ is a subspace of $V$, hence
finite-dimensional, and $\|\cdot\|_\infty$ is a genuine norm on it.  By the
triangle inequality, distinct columns satisfy
\[
 \|A_j-A_k\|_\infty\ \ge\ \|C_j-C_k\|_\infty-2\varepsilon
 \ =\ 1-t-2\varepsilon\ >\ 0 .
\]
The compactness input is the standard fact that closed bounded subsets of
a finite-dimensional real normed space are compact. The use of
$W=V\cap\ell^\infty$ is essential: the individual separated row factors
need not be bounded.
The columns $A_j$, $j\ge0$, are therefore infinitely many points of $W$, all of
norm at most $1+\varepsilon$ and pairwise separated by a fixed positive
distance.  A bounded subset of a finite-dimensional normed space is totally
bounded, so it contains no infinite uniformly separated family.  This
contradiction gives $\varepsilon\ge(1-t)/2$ for every $A$ of finite separated
rank.

For sharpness take $A(i,j)=(1+t)/2$, which has separated rank one; every entry
of $C$ is at distance exactly $(1-t)/2$ from it.
\end{proof}

The compactness step is standard, and the statement is recorded here for this
kernel without a claim of technique.  Three qualifications fix its
reach.  It is a statement about the normalised carry
matrix~\eqref{long269:eq:carry-matrix}: the original kernel carries positive row and
column factors that decay, so the conclusion transfers to a weighted uniform
norm relative to those factors and not to the unweighted sup norm of
$\Kernel$.  It bounds approximation of the whole infinite matrix, and on any
finite range a separated approximant of small rank exists.  And it says nothing
about approximating the scalar sum, so it yields no irrationality conclusion.
What it does add to Theorem~\ref{long269:res:infinite-rank} is robustness: no finite
separated model of the carry, however chosen, gets uniformly closer than half a
carry jump.  The threshold is exact in both directions, since rank one attains
it.

The original kernel behaves differently in the summation norm because its row
and column factors decay.  Let $K^{(N)}(i,j,k)=\Kernel(i,j,k)$ for $i<N$ and
$K^{(N)}(i,j,k)=0$ otherwise.  This is a sum of at most $N$ terms separated
between $i$ and $(j,k)$.  Since
$\hgt(x)>x^3/(pqr)$ for $x>0$, geometric summation gives
\begin{equation}\label{long269:eq:l1-finite-rank-approximation}
 \sum_{i,j,k\ge0}|\Kernel(i,j,k)-K^{(N)}(i,j,k)|
 \le
 \frac{pqr\,p^{-3N}}
 {(1-p^{-3})(1-q^{-3})(1-r^{-3})}.
\end{equation}
Hence the original summable kernel has finite separated-rank approximants in
$\ell^1$, even though its normalised carry matrix has the sharp uniform
barrier of Theorem~\ref{long269:res:uniform-rank}.  This norm distinction is
why neither statement decides the arithmetic nature of the scalar sum.

\section{Dyadic blocks and the literal infinite tail}\label{long269:sec:blocks}

For the rest of the note set $P=\{2,3,5\}$ and $S=\mathcal R_P$, and write
$\hgt_a=\hgt(2^{a})$ and $h_a=\hgt_a/2$, so that $h_0=1/2$ while $h_a$ is a
positive integer for $a\ge1$.  Define the dyadic shell mass, its tail and its
normalisation by
\begin{equation}\label{long269:eq:actual-tail}
 s_a=\sum_{\substack{i,j,k\ge0\\ 2^{a}\le2^{i}3^{j}5^{k}<2^{a+1}}}
       \frac1{\hgt(2^{i}3^{j}5^{k})},
 \qquad
 U_a=\sum_{j\ge a}s_j,
 \qquad
 X_a=h_aU_a .
\end{equation}

Compress the jump word of Section~\ref{long269:sec:lcm} into the blocks cut out by
consecutive powers of two: block $a$ starts just after $2^{a}$, contains every
pure $3$- or $5$-power strictly between $2^{a}$ and $2^{a+1}$, and ends with
the jump at $2^{a+1}$.  A channel cannot occur twice inside one block, because
two powers of the same base $b\ge2$ inside an interval of ratio $2\le b$ must
coincide; this is the
checked internal-power uniqueness lemma for the
internal-power predicate.  Let $I_a$ list the internal jumps $(p,e)$
with $p\in\{3,5\}$ and $2^{a}<p^{e}<2^{a+1}$, in increasing order.

\begin{proposition}[the dyadic block alphabet]\label{long269:res:dyadic-alphabet}
For every $a$,
\begin{equation}\label{long269:eq:dyadic-alphabet}
 b_a=\frac{\hgt_{a+1}}{\hgt_a}=2\prod_{(p,e)\in I_a}p\in\{2,6,10,30\},
 \qquad\text{so}\qquad 2\le b_a\le30 .
\end{equation}
\end{proposition}

\begin{proof}
Internal-power uniqueness leaves two independent yes-or-no choices, one for the
$3$-channel and one for the $5$-channel.  Multiplying the terminal factor $2$
by the selected channel factors gives exactly the four displayed cases.
\end{proof}



For $p\in P$ with $q,r$ the other two primes, put
\[
 A_p(e)=\#\{(i,j)\in\N^2:q^{i}r^{j}<p^{e}\},\qquad
 C_p(e)=\sum_{u=1}^{e}A_p(u),\qquad C_p(0)=0 ,
\]
and for $(p,e)\in I_a$ let $\sigma_a(p,e)$ be the product of the channels of
the later internal jumps of block $a$.  The \emph{ordered block digit} is
\begin{equation}\label{long269:eq:actual-digit}
 m_a=A_2(a+1)+\sum_{(p,e)\in I_a}(p-1)\,\sigma_a(p,e)\,
        \bigl(C_p(e)-C_2(a)\bigr)\quad(a\ge1),
 \qquad m_0=1 ,
\end{equation}
the integer implemented by the pinned checker and the
checked ordered digit.  The first four pairs $(b_a,m_a)$ from $a=1$ are
$(6,4)$, $(10,7)$, $(6,7)$, $(30,65)$.

\begin{theorem}[the literal shell recurrence]\label{long269:res:actual-orbit}
The shell masses are summable and $S=\sum_{a\ge0}s_a$.  For every $a\ge0$,
\begin{equation}\label{long269:eq:shell-digit-identity}
 m_a=h_{a+1}s_a\in\N_{>0},
 \qquad
 X_{a+1}=b_aX_a-m_a,
 \qquad
 X_a=\sum_{j\ge a}\frac{m_j}{b_ab_{a+1}\cdots b_j} .
\end{equation}
\end{theorem}

\begin{proof}
Every exponent of $3$ or $5$ in the $a$th shell is at most $a$, and for each
such pair Lemma~\ref{long269:res:short} leaves at most one exponent of $2$ placing the
point in $[2^{a},2^{a+1})$.  Each height is at least $2^{a}$ and the shell
contains $2^{a}$, so $0<s_a\le(a+1)^{2}2^{-a}$.  The majorant is summable,
which justifies every tail splitting below, and unique prime factorisation
identifies $\sum_a s_a$ with the original repeated series.

Write the internal jumps as $t_\ell=p_\ell^{e_\ell}$, $1\le\ell\le v$, and set
$t_0=2^{a}$, $t_{v+1}=2^{a+1}$.  Let $N(t)$ count smooth positive integers
strictly below $t$ and put $n_\ell=N(t_\ell)$.  On $[t_\ell,t_{\ell+1})$ the
height is $\hgt_a\prod_{u\le\ell}p_u$, and the terminal jump has factor two, so
\[
 h_{a+1}s_a
 =\sum_{\ell=0}^{v}(n_{\ell+1}-n_\ell)\prod_{u>\ell}p_u
 =n_{v+1}-n_0+\sum_{\ell=1}^{v}(p_\ell-1)
        \Bigl(\prod_{u>\ell}p_u\Bigr)(n_\ell-n_0),
\]
the second equality being finite summation by parts.  Counting by the exponent
of $p$ gives $N(p^{e})=C_p(e)$, so the right-hand side is exactly
\eqref{long269:eq:actual-digit}; at $a=0$ the shell is the single point $1$, giving
$m_0=1$.  Positivity follows from $s_a>0$.  Splitting $U_a=s_a+U_{a+1}$ and
multiplying by $h_a$ gives the recurrence, and
$b_a\cdots b_j=\hgt_{j+1}/\hgt_a$ with $m_j=h_{j+1}s_j$ gives
$m_j/(b_a\cdots b_j)=h_as_j$, whose summation is the last identity.
\end{proof}



\begin{lemma}[the literal weighted lattice triangle]
\label{long269:res:literal-triangle}
Put $\lambda_3=\log_2 3$, $\lambda_5=\log_2 5$ and
$\theta_p=1/\lambda_p$ for $p=3,5$.
For $j,k\ge0$ write $w_{j,k}=j\lambda_3+k\lambda_5$ and
$t_{j,k}=\{w_{j,k}\}$.  With $P_a=\hgt(2^a)$, the actual shell numerator is
\[
 m_a=\sum_{\substack{j,k\ge0\\w_{j,k}<a+1}}
 3^{\lfloor(a+1)\theta_3\rfloor-\lfloor(a+t_{j,k})\theta_3\rfloor}
 5^{\lfloor(a+1)\theta_5\rfloor-\lfloor(a+t_{j,k})\theta_5\rfloor}.
\]
Every summand is in $\{1,3,5,15\}$.  In particular, for $a\ge0$,
\[
 \bigl(\lfloor a/(2\lambda_3)\rfloor+1\bigr)
 \bigl(\lfloor a/(2\lambda_5)\rfloor+1\bigr)
 \le m_a\le15(a+1)^2.
\]
Thus $m_a=\Theta((a+1)^2)$ and the numerator sequence is unbounded.
These are integer forcing numerators, not positional digits restricted to
$\{0,\ldots,b_a-1\}$.
\end{lemma}
\begin{proof}
For each pair with $w_{j,k}<a+1$, exactly one exponent
$i=a-\lfloor w_{j,k}\rfloor\ge0$ puts
$2^i3^j5^k=2^{a+t_{j,k}}$ in $[2^a,2^{a+1})$.
This is a bijection onto the literal shell.  Substituting in
$P_{a+1}/(2\hgt(2^i3^j5^k))$ cancels the two-power and gives the displayed
weight.  Each remaining floor increment is zero or one since
$0<\theta_p<1$.  Finally $\lambda_3,\lambda_5>1$ implies
$j,k\le a$ for every summation pair, proving the upper bound.
For the lower bound, restrict to
$0\le j\le\lfloor a/(2\lambda_3)\rfloor$ and
$0\le k\le\lfloor a/(2\lambda_5)\rfloor$.
Then $w_{j,k}\le a$ and each weight is at least one.
\end{proof}
This is an exact reindexing, not a polynomial formula for $m_a$.
It isolates the moving triangular boundary as well as the two floor
weights.  The displayed proof and the finite checks accompanying this
revision are not claimed as a new Lean declaration.

Since $h_0=1/2$, the last identity at $a=0$ writes $S/2=X_0$ as the Cantor
series $\sum_{j\ge0}m_j/(b_0\cdots b_j)$, and $X_a$ is its normalised tail.
Such tails, with their integer carry recurrence, are the classical tool for
the rationality of Cantor series.  Erd\H{o}s and Straus characterise
rationality by integer carries
\cite[Theorem~2.1 and (2.4), pp.~85--86]{erdosstraus1974}, and Han\v{c}l and
Tijdeman give a practical form of that criterion through tails $R_n$ with
$a_nR_n=b_n+R_{n+1}$, which is the recurrence above in their notation
\cite[\S\S2--3 and Theorem~3.1, pp.~372--375]{hancltijdeman2004}.
Koutsoukou-Argyraki and Li formalised the Erd\H{o}s--Straus criteria in
Isabelle/HOL~\cite{afperdosstraus2020}.  Both criteria assume that the $n$th
numerator is $o(a_{n-1}a_n)$, where $a_n$ is the $n$th base.  Here $m_a\ge1$
and $b_{a-1}b_a\le900$, so neither criterion applies.  The passage from
rationality to integral tails is the standard denominator clearing, proved for
this series in Section~\ref{long269:sec:actual-orbit}; the work specific to
this series is the literal forcing with its multiplicities, the strict
endpoint normalisation and the explicit bounds.
For polynomial Cantor data, Han\v{c}l--Tijdeman
\cite[Theorems~2.2 and 3.1]{hancltijdeman2008} give a polynomial
cancellation criterion and a division mechanism; their nonconstant
polynomial radix is not our bounded $(b_a)$.
Their separate Theorem~4.2 permits general integer radices, but assumes an
exact finite-product decomposition with an integer transformed numerator
$o(b_a)$.  Here boundedness would force that transformed numerator to vanish
eventually.  Such a decomposition has not been proved for the literal
$m_a$; the bound $m_a=O(a^2)$ does not produce one, nor does it force an
integer carry to be a polynomial in the floor coordinates.


\begin{proposition}[integral state or cofinal separation]\label{long269:res:actual-dichotomy}
For every integer $B\ge1$, either $BX_a\in\Z$ for some $a\ge0$ and every
later $a$, or for every $a_0$ there is $a\ge a_0$ with
$|BX_a-z|\ge1/31$ for every $z\in\Z$.
\end{proposition}

\begin{proof}
Set $Y_a=BX_a$, so $Y_{a+1}=b_aY_a-Bm_a$.  If cofinal separation fails,
there are $A$ and integers $z_a$ with $e_a=Y_a-z_a$ and $|e_a|<1/31$
for all $a\ge A$.  The recurrence gives
$z_{a+1}-b_az_a+Bm_a=b_ae_a-e_{a+1}$.  The left side is an integer and the right
side has absolute value below $(30+1)/31=1$, so both vanish and
$e_{a+1}=b_ae_a$.  Hence $|e_{A+k}|\ge2^{k}|e_A|$ for every $k$ while
$|e_{A+k}|<1/31$, which forces $e_A=0$.  Once $Y_A$ is integral, the
integer recurrence makes every later $Y_a$ integral.
\end{proof}

The supplied formal source records the $B=1$ instance.  Multiplying the
forcing by the integer $B$ gives the ordinary proof of the scaled statement
above, without requiring a new infinite-tail identity.  Neither branch is
excluded.  For comparison, Erd\H{o}s--Taylor
\cite[Theorem~1, p.~600]{erdostaylor1957} prove a countability statement for
increasing integer sequences with bounded successive ratios; Fan
\cite[Lemma~3.1]{fan2026strongly} gives the unbounded-sequence version.
For $g_a=Bh_a$ our integer-ratio argument is stronger: convergence of
$\operatorname{dist}(g_a\xi,\Z)$ to zero forces eventual integrality.
Applying that statement to the particular $\xi=S$ still requires excluding
the integral branch.

\section{A quadratic tail bound and denominator cancellation}
\label{long269:sec:actual-orbit}

Put
\[
 n_a=a+\lfloor\log_3(2^{a})\rfloor+\lfloor\log_5(2^{a})\rfloor,
 \qquad
 Q(n)=\frac{n^{2}+8n+18}{9},
\]
so $n_a$ is the sum of the three height exponents at $2^{a}$.

\begin{theorem}[quadratic bound for the actual tail]\label{long269:res:actual-tail-bound}
For every $a\ge0$, $0<X_a\le Q(n_a)$.
\end{theorem}

\begin{proof}
Partition the smooth integers $x\ge2^{a}$ by their height vector
\[
 (A,B,C)=(\lfloor\log_2x\rfloor,\lfloor\log_3x\rfloor,\lfloor\log_5x\rfloor).
\]
The cell of a vector is the interval $[\lo,\hi)$ with
$\lo=\max(2^{A},3^{B},5^{C})$ and $\hi=\min(2^{A+1},3^{B+1},5^{C+1})$, so
$\hi\le2^{A+1}\le2\lo$ and, by Lemma~\ref{long269:res:short}, fixing the exponents of
$3$ and $5$ leaves at most one exponent of $2$.  Since $A\ge B\ge C$, the cell
contains at most
\[
 (B+1)(C+1)\le\frac{(A+B+C+3)^{2}}{9}
\]
smooth points: writing $u=B+1\ge v=C+1$ and using $A+1\ge u$, the difference
$(2u+v)^{2}-9uv=(u-v)(4u-v)$ is nonnegative.  Unique factorisation identifies
these exponent triples with distinct smooth integers.

As the cutoff increases all three height exponents are nondecreasing, so two
nonempty cells with the same exponent sum are the same cell, and there is at
most one nonempty cell of each rank.  Every cell above $2^{a}$ has rank at
least $n_a$, and a cell of rank $n_a+k$ has height at least $\hgt_a2^{k}$,
since each of the $k$ extra prime factors is at least $2$.  Nonnegative
summation over ranks, allowing empty ones, gives
\[
 X_a\le\frac1{18}\sum_{k\ge0}\frac{(n_a+k+3)^{2}}{2^{k}}
      =\frac{n_a^{2}+8n_a+18}{9},
\]
the evaluation using the geometric moments $\sum2^{-k}=2$, $\sum k2^{-k}=2$
and $\sum k^{2}2^{-k}=6$.  Positivity follows from the shell at $2^{a}$.
\end{proof}

The rank here is the sum of the three fixed height exponents, not a
varying-smoothness density parameter. The estimate uses only unique
factorisation, the short-cell counting bound and geometric moments; it does
not invoke a Dickman--Hildebrand asymptotic or a Hecke--Mahler theorem.

This is the analytic content of the section: each additional height rank costs
a geometric factor while its multiplicity grows only quadratically.  The finite
projection and the sorted quadratic inequality are the checked ingredients of
Proposition~\ref{long269:res:drop}; grouping the actual infinite tail by height cells
and summing the majorant is the argument above.

The bound is used at the endpoint of a window, where the natural index is the
positive prime-power jump count strictly below the cutoff.  Put
\begin{equation}\label{long269:eq:endpoint-index}
 j_a=\#\{p^{e}<2^{a}:p\in\{2,3,5\},\ e\ge1\}=n_a-1\qquad(a\ge1),
\end{equation}
the equality holding for $a\ge1$ because the powers of $2$ below $2^{a}$ number
$a-1$ while the powers of $3$ and of $5$ below $2^{a}$ number
$\lfloor\log_3 2^{a}\rfloor$ and $\lfloor\log_5 2^{a}\rfloor$; at $a=0$ the
count is $j_0=0$ and $n_0-1=-1$.  Substituting $n_a=j_a+1$ into $Q$ gives the
integer bound used throughout the rest of the note:
\begin{equation}\label{long269:eq:actual-bound}
 K(B,a)=\bigl\lfloor B\,Q(n_a)\bigr\rfloor,
 \qquad
 K(B,a)=\left\lfloor\frac{B(j_a^{2}+10j_a+27)}9\right\rfloor\quad(a\ge1).
\end{equation}
The cutoff in \eqref{long269:eq:endpoint-index} is $2^{a}$ and it is strict.
The symbol $K(B,a)$ is an integer cap for an \emph{integral} quantity
$BX_a$: from $BX_a\le BQ(n_a)$ one may take the floor only after integrality
has been established.  We do not assert $BX_a\le K(B,a)$ for arbitrary
real tails.


\begin{lemma}[finite denominator clearing]\label{long269:res:all-scale-lattice}
For all integers $0\le u\le b$ the window mass $h_b\sum_{a=u}^{b-1}s_a$ is a
natural number.  If $S=N/D$ with $N\in\Z$ and $D\in\N_{>0}$, then $DX_a\in\Z$
for every $a\ge1$, and there are indices $1\le i<j\le D+1$ for which
$X_i-X_j\in\Z$.
\end{lemma}

\begin{proof}
An empty window has mass zero.  Otherwise $b\ge1$, and every integer $x<2^{b}$
has $2$-height exponent at most $b-1$ while its other height exponents are at
most those at $2^{b}$, so $\hgt(x)\mid\hgt_b/2=h_b$ and every term of the
finite window clears at $h_b$.  Writing $v_a=h_a\sum_{u<a}s_u\in\N$ gives
$DX_a=h_aN-Dv_a\in\Z$.  Among $D+1$ of these integers two share a residue
modulo $D$, and the corresponding states differ by an integer.
\end{proof}

The strict upper endpoint is what permits division by two.  A cutoff
including $2^b$ would not clear its term at the normaliser $h_b=P_b/2$.

\begin{theorem}[positive reduced carries from rationality]
\label{long269:res:actual-cancellation}\label{long269:res:lead-carry-bridge}
\label{long269:res:actual-carry-bound}\label{long269:res:denominator-reduction}
Suppose $S=N/D$ with $N\in\Z$, $D\in\N_{>0}$, and write
\[
 D=2^{u}3^{v}5^{w}B,\qquad u,v,w\in\N,\quad B\in\N_{>0},\quad\gcd(B,30)=1,
 \qquad a_D=u+1+2v+3w .
\]
Then for every $a\ge a_D$ the number $z_a=BX_a$ is a positive integer and
\[
 z_{a+1}=b_az_a-Bm_a,\qquad 1\le z_a\le K(B,a)\le90B(a+1)^{2} .
\]
\end{theorem}

\begin{proof}
Let $M=2^{u}3^{v}5^{w}$.  For $a\ge a_D$ we have $2^{a}\ge2^{u+1}$,
$2^{a}\ge3^{v}$ and $2^{a}\ge5^{w}$, using $3<2^{2}$ and $5<2^{3}$; hence
$M\mid h_a$.  In the clearing identity $DX_a=h_aN-Dv_a$ of
Lemma~\ref{long269:res:all-scale-lattice} both terms on the right are divisible by $M$,
so dividing by $M$ shows that $BX_a$ is an integer.  Positivity and the
recurrence come from~\eqref{long269:eq:shell-digit-identity}, and the upper bound is
Theorem~\ref{long269:res:actual-tail-bound} with the floor taken, since $z_a$ is an
integer below $BQ(n_a)$.  Finally $n_a\le3a$, so
$Q(n_a)\le a^{2}+\tfrac83a+2\le90(a+1)^{2}$.
\end{proof}

The smooth part $M$ affects the clearing onset; the surviving carry bound
depends on $B$.  The displayed $a_D$ is a convenient sufficient onset.
The next proposition gives the exact onset when $N/D$ is in lowest terms.

\begin{proposition}[exact denominators and minimal clearing]
\label{long269:res:exact-denominator}
Suppose $S=N/(MB)$ is in lowest terms, with $M=2^u3^v5^w$ and
$\gcd(B,30)=1$.  For every $a\ge1$,
\[
 \operatorname{den}(X_a)=\frac{MB}{\gcd(M,h_a)},\qquad
 \operatorname{den}(BX_a)=\frac{M}{\gcd(M,h_a)}.
\]
Hence the first integral reduced tail occurs at
\[
 a_* =\min\{a\ge1:2^a\ge\max(2^{u+1},3^v,5^w)\}\le a_D,
\]
and $BX_a$ is integral exactly for $a\ge a_*$.  This onset is computable by
integer powers; it is not an estimate obtained by rounding logarithms.
\end{proposition}
\begin{proof}
Write $X_a=h_aN/(MB)-v_a$, with $v_a\in\mathbb Z$ as in
Lemma~\ref{long269:res:all-scale-lattice}.  Since $N$ is coprime to $MB$
and $h_a$ is supported on $\{2,3,5\}$, reduction gives both denominators.
Now $M\mid h_a$ means
$a-1\ge u$, $\lfloor\log_3 2^a\rfloor\ge v$ and
$\lfloor\log_5 2^a\rfloor\ge w$, precisely the three integer inequalities.
All three persist when $a$ increases, giving the first and every later
integral reduced tail.  The earlier bounds $3<4$ and $5<8$ give $a_*\le a_D$.
\end{proof}
This is an ordinary refinement of the supplied sufficient-onset bridge,
not a claim of a newly checked Lean theorem.

The integral branch has one further exact property, recorded because it
constrains any attempt to build a surviving seed by hand.

\begin{proposition}[upward closure and rigidity]\label{long269:res:pinning}
For every $a$, $X_a=(m_a+X_{a+1})/b_a>0$, and if $X_a\in\Z$ then $X_n\in\Z$ for
every $n\ge a$.  Moreover, fix $A$, a positive width function $w$ with
$w(A+k)/8^{k}\to0$, and a real orbit $(y_n)_{n\ge A}$ satisfying
$y_{n+1}=b_ny_n-m_n$.  If $y_n$ and $X_n$ both lie in
$(m_n/b_n,\;m_n/b_n+w(n)]$ for every $n\ge A$, then $y_A=X_A$.
\end{proposition}

\begin{proof}
The identity is the recurrence solved for $X_a$, and positivity holds because
every shell contains its dyadic left endpoint.  Integer coefficients give
upward closure.  For the last assertion,
$y_{A+k}-X_{A+k}=(P_{A+k}/P_A)(y_A-X_A)$.
The height quotient satisfies $P_{A+k}/P_A>8^k/15$ for $k\ge1$
(Proposition~\ref{long269:res:window-growth}), whereas the common interval
bounds the absolute difference by $w(A+k)$.  Thus
$|y_A-X_A|<15w(A+k)/8^k\to0$.
\end{proof}

\subsection{The jump-constrained improvement}
\begin{proposition}[a smaller bound for the actual tail]
\label{long269:res:jump-constrained-bound}
For every $a\ge0$,
\[
 0<X_a\le\widetilde Q(n_a)<Q(n_a),\qquad
 \widetilde Q(n)=\frac{1210n^2+9130n+18847}{11979}.
\]
\end{proposition}
\begin{proof}
For a point $x\ge2^a$, let $d_p$ be the increase in its $p$-height exponent
from the boundary $2^a$, and put $k=d_2+d_3+d_5$.
Writing $B=\lfloor\log_3 2^a\rfloor$, we have
$x<3^{B+d_3+1}\le2^a4^{d_3+1}$, hence $d_2\le2d_3+1$.
Thus $j=\lfloor(k+1)/3\rfloor\le d_3+d_5$ and
\[
 \frac{\hgt(x)}{P_a}\ge2^{k-j}3^j=:d(k).
\]
There is at most one height cell of each rank, with at most
$(n_a+k+3)^2/9$ supported points by the earlier projection bound.
Consequently
\[
 X_a\le\sum_{k\ge0}\frac{(n_a+k+3)^2}{18d(k)}.
\]
Now $d(3m)=12^m$, $d(3m+1)=2\cdot12^m$ and
$d(3m+2)=6\cdot12^m$.  Grouping in threes and summing the quadratic
geometric series gives $\widetilde Q(n_a)$.
Finally $11979\bigl(Q(n)-\widetilde Q(n)\bigr)
=121n^2+1518n+5111>0$ for $n\ge0$.
\end{proof}
The supplied public Lean at commit
\nolinkurl{3d6d938d696fed0fb71dd55115a18a73738ff223} proves this for the
\emph{actual} normalised tail, not only for an abstract jump model:
\texttt{actual\_sharp\_tail\_bound} in
\texttt{ActualSharpTailMajorantR10.lean}, namespace
\texttt{ErdosProblems.Erdos269.PaperR10}
(line 339 in the attached file).  The smaller integer cap
$\widetilde K(B,a)=\lfloor B\widetilde Q(n_a)\rfloor$ is therefore
available to the rational direction.  We retain $K$ in the main equivalence
and old finite certificates to avoid silently changing the recorded tests.
No optimality claim is made for either quadratic bound.  The rational
implication also works for any $G\ge\widetilde K$; together with
$G(B,a)=o(8^a)$, the same proof gives an equivalence on this enlarged band.
The named theorem below deliberately retains $K$ so that its formal-source
interface and the archived certificates are unchanged.

\section{Residue windows and the equivalence band}\label{long269:sec:escape}

The bases and accumulated forcing in this section are integer sequences;
the actual tails remain real unless rationality has been assumed.  For
$\ell\ge1$ and $h\ge0$, set
\[
 W_{\ell,0}=1,\quad F_{\ell,0}=0,\qquad
 W_{\ell,h+1}=b_{\ell+h}W_{\ell,h},\quad
 F_{\ell,h+1}=b_{\ell+h}F_{\ell,h}+m_{\ell+h}.
\]
Induction gives
\begin{equation}\label{long269:eq:window-identity}
 X_{\ell+h}=W_{\ell,h}X_\ell-F_{\ell,h},\qquad
 z_{\ell+h}=W_{\ell,h}z_\ell-BF_{\ell,h}
\end{equation}
for any orbit satisfying $z_{n+1}=b_nz_n-Bm_n$.
For integers $C\ge1,N$, use
$\lpr_C(N)=1+((N-1)\bmod C)\in\{1,\ldots,C\}$.
In particular $\lpr_C(0)=C$.  Positivity of this representative is what
allows comparison with a positive integral carry.

\begin{proposition}[the finite residue contradiction]\label{long269:res:consumer}
Let $C>0$ and let $c$ be an integer with $0<c$ and $|c|\le K$.  If
$c\equiv N\pmod C$ and $K<\lpr_C(N)$, then the hypotheses are contradictory.
\end{proposition}

\begin{proof}
The canonical representative lies in $\{1,\ldots,C\}$, and the inequalities put
$c$ strictly between $0$ and $C$.  If $N\equiv0\pmod C$ its representative is
$C$ while $c\bmod C=c\ne0$.  Otherwise $c$ and $\lpr_C(N)$ are each their own
residue and congruence makes them equal, contradicting $|c|\le K<\lpr_C(N)$.
\end{proof}



For a bound $G:\N_{>0}\times\N\to\N$ let $\Esc(G)$ be the statement
\begin{equation}\label{long269:eq:actual-escape}
 \begin{gathered}
 \text{for every }B\ge1\text{ with }\gcd(B,30)=1\text{ and every }a_0\ge1,\\
 \text{there are }\ell\ge a_0\text{ and }h\ge1\text{ with }
 \lpr_{W_{\ell,h}}(-BF_{\ell,h})>G(B,\ell+h) .
 \end{gathered}
\end{equation}
The window may depend on both $B$ and $a_0$.  All quantities in the
inequality are finite integers.  The quantifiers over denominators and
prescribed onsets are nevertheless unbounded; a search over a finite
rectangle of $(B,\ell)$ does not verify them.  Every $b_a>0$, so
$W_{\ell,h}>0$ is automatic.

\begin{theorem}[the equivalence band]\label{long269:res:actual-escape-endpoint}
\label{long269:res:lead-escape-equivalence}\label{long269:res:windowconsumer}
Let $G:\N_{>0}\times\N\to\N$ satisfy $K(B,a)\le G(B,a)$ for all $B$ and $a$,
and $G(B,a)/8^{a}\to0$ as $a\to\infty$ for each fixed $B$.  Then
\[
 \Esc(G)\quad\Longleftrightarrow\quad S\notin\Q .
\]
Both $K$ of~\eqref{long269:eq:actual-bound} and $K_0(B,a)=90B(a+1)^{2}$ lie in this
band, so $\Esc(K)$, $\Esc(K_0)$ and irrationality of $S$ are mutually
equivalent.  The lower half of the band cannot be dropped: $\Esc(0)$ holds
automatically, since every least positive residue is at least $1$.
\end{theorem}

\begin{proof}
Suppose $\Esc(G)$ and suppose $S=N/D$ were rational.
Theorem~\ref{long269:res:actual-cancellation} supplies $B\ge1$ coprime to $30$, an
onset $a_D$, and positive integers $z_a=BX_a\le K(B,a)\le G(B,a)$ for
$a\ge a_D$ satisfying the cleared recurrence.  Apply~\eqref{long269:eq:actual-escape}
with $a_0=a_D$ to obtain a window $(\ell,h)$ with $\ell\ge a_D$.
By~\eqref{long269:eq:window-identity}, $z_{\ell+h}\equiv-BF_{\ell,h}$ modulo
$W_{\ell,h}$, while $0<z_{\ell+h}\le G(B,\ell+h)<\lpr_{W_{\ell,h}}(-BF_{\ell,h})$.
Proposition~\ref{long269:res:consumer} is the contradiction, so $S$ is irrational.

Conversely suppose $S\notin\Q$, and fix $B\ge1$ coprime to $30$ and
$a_0\ge1$.  Set $\ell=a_0$ and
$\delta=\lceil BX_\ell\rceil-BX_\ell\in(0,1)$; the prefix identity makes
$BX_\ell$ irrational.  For all sufficiently large $h$,
$0<\delta+BX_{\ell+h}/W_{\ell,h}<1$, because
$X_{\ell+h}=O((\ell+h+1)^2)$ and $W_{\ell,h}>8^h/15$.
The integer $\lceil BX_\ell\rceil W_{\ell,h}-BF_{\ell,h}$ is therefore
exactly
\[
 \lpr_{W_{\ell,h}}(-BF_{\ell,h})
 =\delta W_{\ell,h}+BX_{\ell+h}.
\]
Its first term eventually exceeds $G(B,\ell+h)$, since
$G(B,a)=o(8^a)$.  Thus every sufficiently long window from this fixed
start escapes, which is stronger than the required existence.

For the two named bounds, $K\le K_0$ by
Theorem~\ref{long269:res:actual-cancellation}, both are quadratic in $a$ for fixed $B$,
and $K\le K$ trivially.  Finally $\lpr_C(N)\ge1>0$ always, so $\Esc(0)$ holds
whatever $S$ is, while the required domination $K\le0$ fails.
\end{proof}

The supplied source records the fixed-cap equivalence and the larger
$o(8^a)$ band with domination of $K$, in separate modules listed in the
source concordance.  The exact fixed-start residue identity used above is
an ordinary refinement of the proof.  No fresh formal-system build was
performed for this revision.

\begin{proposition}[the fixed-start residue limit]
\label{long269:res:residue-limit}
For fixed integers $B,\ell\ge1$, write
$R_h=\lpr_{W_{\ell,h}}(-BF_{\ell,h})$.  Then
\[
 \lim_{h\to\infty}\frac{R_h}{W_{\ell,h}}=
 \begin{cases}
 \lceil BX_\ell\rceil-BX_\ell,&BX_\ell\notin\mathbb Z,\\
 0,&BX_\ell\in\mathbb Z.
 \end{cases}
\]
In the nonintegral case,
$R_h=(\lceil BX_\ell\rceil-BX_\ell)W_{\ell,h}+BX_{\ell+h}$
for all sufficiently large $h$.  In the integral case,
$R_h=BX_{\ell+h}$ for all sufficiently large $h$.
\end{proposition}
\begin{proof}
The preceding argument uses only nonintegrality of $BX_\ell$, not
irrationality of $S$, and proves the first identity.  If $BX_\ell$ is
integral, the recurrence makes $BX_{\ell+h}$ a positive integer congruent
to $-BF_{\ell,h}$.  Eventually it is less than $W_{\ell,h}$, so it is the
least positive representative.  In both cases
$BX_{\ell+h}/W_{\ell,h}\to0$.
\end{proof}
For a rational value, a start before the clearing onset can therefore
admit arbitrarily long escaping windows.  This is why the rational
contradiction must choose a start after the onset; one early start alone
does not rule out rationality.

Theorem~\ref{long269:res:actual-escape-endpoint} makes the producer a
restatement of Erd\H{o}s~\#269 for $P=\{2,3,5\}$ throughout the band, so work
on it is work on the target.  The rational direction needs a bound valid for
every actual positive carry.  The bound $K$ is one such bound, and a smaller
bound proved valid for the actual carries could replace it; the zero bound is
invalid, and $\Esc(0)$ holds vacuously.  Enlarging the bound to any $G\ge K$
with $G(B,a)=o(8^{a})$, in particular to any subexponential bound above $K$,
preserves the equivalence.  Equivalence preserves truth and settles nothing
about difficulty, so the reformulation may still be the easier representation
to attack.

Two exact countermodels rule out further shortcuts.  For $(W,F,B)=(6,4,1)$ we
have $\lpr_6(-4)=2$, so the canonical residue need not be coprime to the
accumulated base.  For $(W,F,B)=(60,47,37)$ we have
$\lpr_{60}(-37\cdot47)=1$, so a fixed window has no denominator-independent
lower bound on that residue.  The window may therefore depend genuinely on $B$.

\subsection{How fast the window base grows}

The crude estimate $W_{\ell,h}\ge2^{h}$ proves the smaller $o(2^a)$ band.
The exact height formula proves the $o(8^a)$ band above and sets the scale
of any finite search.

\begin{proposition}[window-growth law]\label{long269:res:window-growth}
Put $\theta_3=\log_32$ and $\theta_5=\log_52$.  For all $\ell\ge0$ and
$h\ge1$,
\[
 W_{\ell,h}=2^{h}\,
 3^{\lfloor(\ell+h)\theta_3\rfloor-\lfloor\ell\theta_3\rfloor}\,
 5^{\lfloor(\ell+h)\theta_5\rfloor-\lfloor\ell\theta_5\rfloor},
 \qquad
 \frac{8^{h}}{15}<W_{\ell,h}<15\cdot8^{h} .
\]
\end{proposition}

\begin{proof}
Telescoping \eqref{long269:eq:dyadic-alphabet} gives
$W_{\ell,h}=\hgt_{\ell+h}/\hgt_\ell$, and the displayed formula is that
quotient written out.  Each floor difference differs from $h\theta_p$ by less
than one, and $3^{\theta_3}=5^{\theta_5}=2$, so the $3$-factor lies strictly
between $2^{h}/3$ and $3\cdot2^{h}$ and the $5$-factor strictly between
$2^{h}/5$ and $5\cdot2^{h}$.  Multiplying the three ranges gives the bounds.
\end{proof}

\begin{corollary}[no bounded-length escape at all starts]\label{long269:res:no-bounded-length}
Fix $B\ge1$ coprime to $30$ and $H\ge1$.  Only finitely many starts $\ell$
admit an escaping window of length at most $H$ against the bound $K$.
\end{corollary}

\begin{proof}
A least positive residue never exceeds its modulus, so escape at $(\ell,h)$
requires $K(B,\ell+h)<W_{\ell,h}<15\cdot8^{h}\le15\cdot8^{H}$.  On the other
hand $j_a\ge a-1$, since the powers $2,\ldots,2^{a-1}$ already lie below
$2^{a}$, so $K(B,\ell+h)\ge\lfloor B((\ell-1)^{2}+10(\ell-1)+27)/9\rfloor$,
which tends to infinity with $\ell$.
\end{proof}

\noindent Corollary~\ref{long269:res:no-bounded-length} is the exact reason a finite
scan cannot approach the cofinal quantifier by widening its denominator range
alone.  Rearranging its inequality, an escaping window at start $\ell$ and
endpoint $a=\ell+h$ must satisfy
\[
 3h>\log_2K(B,a)-\log_215,
\]
so the necessary search depth grows like
$\bigl(\log_2B+2\log_2\ell\bigr)/3$.  This is a lower bound on the length that can
possibly work, and it is not an upper bound on the first length that does.
An explicit eventual escape threshold can be obtained from a positive
lower bound for $\delta=\lceil BX_\ell\rceil-BX_\ell$ and for $1-\delta$,
together with the displayed tail and window estimates.  Such information
has not been supplied for the literal series.  A necessary lower bound on
window length is not a prediction of the distribution of first successful
lengths and does not give an upper search bound.

\statusboundary

\section{The finite evidence}\label{long269:sec:evidence}

The three displayed window tuples are directly replayable finite
certificates.  The larger scan and denominator exclusions below are archived
computational reports.  Their full execution records are not supplied here;
in particular the two large exclusions lack their machine-readable witnesses.
They are not theorem inputs, and none supplies the unbounded quantifiers.

\subsection{The dyadic-window scan}

A \emph{certificate} is a finite tuple of integers recording one instance of
the inequality in~\eqref{long269:eq:actual-escape}, so a reader can recheck the modular inequality immediately and reproduce
the shell numerators by exact enumeration.  The integer-only
\href{\sourceurl/scripts/check_erdos269_dyadic_windows.py}{dyadic-window
checker} constructs the ordered pure-power jumps, the block bases and the block
digits from exact multiplicity counts, and is the historical source of the following tuples.  The package supplied
with this revision independently reconstructs and verifies them, with
columns the denominator $B$, the start $\ell$, the length $h$, the endpoint
jump index $j_{\ell+h}$, the window base, the forcing, the residue
$R=\lpr_{W}(-BF)$ and the bound $K$ of~\eqref{long269:eq:actual-bound}.
\[
\begin{array}{c|c|c|c|r|r|r|r}
B&\ell&h&j_{\ell+h}&W&F&R&K\\ \hline
1&1&2&4&60&47&13&9\\
7&1&3&6&360&289&137&95\\
16&1&4&9&10800&8735&640&352
\end{array}
\]
The first row reads as follows.  The window starts at $\ell=1$ and has length
$2$, so $W=b_1b_2=6\cdot10=60$; the accumulated forcing is $F=47$; and
$\lpr_{60}(-47)=13$, since $-47+60=13$, which exceeds
$K(1,3)=\lfloor(16+40+27)/9\rfloor=9$.  The third row lies outside the domain
of~\eqref{long269:eq:actual-escape}, since $\gcd(16,30)=2$, and is displayed to
illustrate the window arithmetic at greater depth.

The archived scan report covers $B\le5000$ coprime to $30$ and
$100\le\ell\le3000$: $3{,}869{,}934$ pairs, with reported first escape
length at most $18$ in a search to length $24$.  This full scan was not
rerun for this revision.  Its histogram, retained in
Section~\ref{long269:long:experiments}, is an archived observation, not a
consequence of the window-growth bound.  A fresh suite supplies all $2496$ witnesses for $1\le B\le97$,
$\gcd(B,30)=1$ and $1\le\ell\le96$, with first escaping lengths from
$1$ to $10$ in a search allowed up to $24$.  A separate integer-only
verifier reconstructs the shells and checks completeness of this rectangle.
Neither suite proves escape for unbounded $B$ or for cofinally many starts;
Corollary~\ref{long269:res:no-bounded-length} rules out covering the latter
quantifier with any fixed maximum length.

\subsection{Two finite denominator exclusions}

\paragraph{Archived window-$128$ exclusion report (witness unavailable).}
\label{long269:res:lead-block-exclusion}
The archived report asserts the following, which is not used as a proved
result in this revision.  At window length $L=128$, launch $a_1=10005$ and $64$ starts, no rational value
of the $\{2,3,5\}$ running-LCM series has reduced denominator $MB$ with $M$ a
$30$-smooth divisor of $2^{10005}3^{6312}5^{4308}$, $\gcd(B,30)=1$ and
$1<B\le B_{\max}$, where $B_{\max}$ is the $106$-digit integer
\[
\begin{aligned}
 B_{\max}={}&1134599670999687767349520845707093359257353022286558739363600235\\
 &016103207564063373270305324172145281971729 ,
\end{aligned}
\]
so that $\log_2B_{\max}=348.9846\ldots$


The exponent triple $(10005,6312,4308)$ is the height exponent triple of
$\hgt(2^{10005})$, so the certificate normalises by the full height at its
launch.  Lemma~\ref{long269:res:all-scale-lattice} normalises by the half height
$h_{10005}$; the certificate's smooth family is accordingly the larger one.  The recorded quantities are the window product with
$\log_2P=386.40993\ldots$, the exclusion index $J=1$ certified from a reduced
basis lying inside the per-start budget, an enclosure of width
$9.674\times10^{-227}$, and a maximum ratio $\max X/W=0.185997$.  The machine-readable certificate and a replay receipt are absent.
Consequently this is a provenance record, not a verified denominator
exclusion available to the reader.  Before restoring a theorem environment,
publish the exact enclosure, reduction basis, integer inequalities, engine
revision and reproducible command.  No Lean declaration is claimed for it.

\paragraph{Archived continued-fraction exclusion report (witness unavailable).}
\label{long269:res:cf-exclusion}
The archived report asserts that the normalised tail $X_1$ has $13{,}109$ certified partial quotients, so it is
not rational with denominator at most $2^{22482}$, about $10^{6768}$.


The stated certification method is a common prefix of the continued fractions of the two
endpoints of an interval provably containing $X_1$; the numbers with a given
prefix of partial quotients form an interval
\cite[Lemma~10.1.2, p.~106]{bosmaCF}, so no approximation heuristic enters,
and the truncation was independently checked to agree with
the direct smooth-number sum as an exact rational. The machine-readable witness and replay receipt are absent; the asserted
exclusion is therefore not a theorem input.  A replay must supply rational
endpoints, a proof that the entire interval contains $X_1$, the common
continued-fraction cylinder, and a rigorous lower bound for denominators of
all rationals in the enclosure.  A count of matching partial quotients
alone is not that denominator certificate.  The recorded statistics are
a largest denominator of $22{,}483$ bits, a largest partial quotient of
$129{,}114$, a mean partial quotient of $23.4133$, observed Gauss--Kuzmin
frequencies $0.4208$, $0.1665$, $0.0917$, $0.0575$, $0.0391$ against the
predicted $0.4150$, $0.1699$, $0.0931$, $0.0589$, $0.0406$, and a L\'evy
constant of $1.18869$ against $\pi^{2}/(12\log2)=1.18657$.  The predicted
values are the almost-everywhere frequencies of Gauss's law and L\'evy's
almost-everywhere constant \cite[pp.~288--289 and footnote~5]{levy1936}.
These statistics describe a finite prefix and do not bear on whether $X_1$ is a
Liouville number, algebraic, or rational with a larger denominator.

The two archived reports concern different finite families: one uses a
lattice at a fixed launch and fixed smooth part, the other a continued-fraction
enclosure at $a=1$.  Neither has a replayable witness in this packet.
Even after verification, each would exclude only its stated denominator
range and neither would settle an instance of the problem.

\section{The remaining arithmetic questions}\label{long269:sec:open}

\subsection{An exact weighted-shift identity for the repeated series}
\label{long269:sec:weighted-shifts}
This subsection concerns $S$, not $\mathcal D_{2,3,5}$.
Write $P_a=\hgt(2^a)$, $\alpha=S/2$ and
$\alpha_t=\sum_{k<t}m_k/P_{k+1}$.  The literal coefficient object is
Lemma~\ref{long269:res:literal-triangle}; its size bound is not a
polynomial representation.

\begin{lemma}[weighted shifts preserve the actual value]
\label{long269:res:weighted-shift-identity}
Fix integers $c_0,\ldots,c_\sigma$, not all zero, independently of the
positive integer shift $r$.  Put
\[
 \gamma_{a,t}=\frac{P_tP_{a+1}}{P_{a+t+1}},\qquad
 D_{r,a}=15\sum_{j=0}^{\sigma}c_j\gamma_{a,jr}m_{a+jr}.
\]
Then $\gamma_{a,t}\in\{1,1/3,1/5,1/15\}$ and $D_{r,a}\in\Z$.
Furthermore, with
\[
 A_r=15\sum_{j=0}^{\sigma}c_jP_{jr},\qquad
 Z_r=15\sum_{j=0}^{\sigma}c_jP_{jr}\alpha_{jr}\in\Z,
\]
we have the absolutely convergent identity
\begin{equation}\label{long269:eq:weighted-shift-identity}
 A_r\alpha-Z_r=\sum_{a\ge0}\frac{D_{r,a}}{P_{a+1}}.
\end{equation}
One may take
$C=225\sum_{j=0}^{\sigma}|c_j|\max(1,j)^2$ in
$|D_{r,a}|\le C(a+r+1)^2$.  If $J$ is the largest index with $c_J\ne0$,
then $A_r\ne0$ whenever
$\sum_{j<J}|c_j|2^{-(J-j)r}<|c_J|$; an empty sum is zero.
\end{lemma}
\begin{proof}
The two-exponents cancel in $\gamma$.  For $p=3,5$, floor addition leaves
exponent zero or $-1$, proving the four-valued assertion and integrality.
For $k<jr$, $P_{k+1}$ divides $P_{jr}$, so $Z_r$ is an integer.
For each $t$, absolute convergence gives
\[
 \sum_{a\ge0}\frac{\gamma_{a,t}m_{a+t}}{P_{a+1}}
 =P_t\sum_{k\ge t}\frac{m_k}{P_{k+1}}
 =P_t(\alpha-\alpha_t).
\]
The finite linear combination proves the identity.
Lemma~\ref{long269:res:literal-triangle}, $0<\gamma\le1$ and
$a+jr+1\le\max(1,j)(a+r+1)$ give the stated constant.
For $j<J$, $P_{jr}/P_{Jr}\le2^{-(J-j)r}$.  Dividing $A_r$ by
$15P_{Jr}$ and bounding the remaining terms proves the nonvanishing
condition, which holds for all sufficiently large $r$.
\end{proof}
The identity was derived in the earlier editorial structural note and is
included here with its complete proof.  It is not claimed as an application
of an existing fixed-base theorem or as a Lean-checked statement.

\subsection{The boundary term in the weighted difference}
\label{long269:sec:boundary-difference}
Let $T_a=\{(j,k)\in\mathbb N^2:w_{j,k}<a+1\}$ and, for $0\le t<1$, set
\[
 \omega_a(t)=\prod_{p\in\{3,5\}}
 p^{\lfloor(a+1)\theta_p\rfloor-\lfloor(a+t)\theta_p\rfloor}.
\]
This weight is defined even for pairs outside $T_a$.  For $\nu\ge0$, put
\[
 \begin{split}
 \kappa_p(a,t,\nu r)
 &=\lfloor(a+t+\nu r)\theta_p\rfloor
   -\lfloor(a+t)\theta_p\rfloor-\lfloor\nu r\theta_p\rfloor,\\
 \chi_\nu(a,t)&=3^{-\kappa_3(a,t,\nu r)}5^{-\kappa_5(a,t,\nu r)}.
 \end{split}
\]
Each $\kappa_p$ is $0$ or $1$.  We call it a crossing only in this precise
floor-addition sense; the moving edge of $T_a$ is a different boundary.

\begin{proposition}[an exact interior-and-strip decomposition]
\label{long269:res:strip-decomposition}
For a fixed operator $c_0,\ldots,c_\sigma$ and $r\ge1$, let
$E_0=T_a$ and $E_s=T_{a+sr}\setminus T_{a+(s-1)r}$ for $1\le s\le\sigma$.
Then
\[
 \frac{D_{r,a}}{15}=
 \sum_{s=0}^{\sigma}\ \sum_{(j,k)\in E_s}
 \omega_a(t_{j,k})\sum_{\nu=s}^{\sigma}c_\nu\chi_\nu(a,t_{j,k}).
\]
For the cubic operator $(1,-3,3,-1)$, if all these crossing bits vanish,
then
\[
 \frac{D_{r,a}}{15}=-M_0+2M_1-M_2,\qquad
 M_s=\sum_{(j,k)\in T_{a+(s+1)r}\setminus T_{a+sr}}\omega_a(t_{j,k}).
\]
Thus absence of floor crossings cancels the common interior, but not
necessarily the three boundary strips.
\end{proposition}
\begin{proof}
Cancellation of the two upper-boundary exponents gives the pointwise identity
\[
 \gamma_{a,\nu r}\omega_{a+\nu r}(t)
   =\omega_a(t)\chi_\nu(a,t).
\]
Insert the triangle formula for each $m_{a+\nu r}$ and group by the first
triangle containing a pair.  A pair in $E_s$ occurs exactly in the terms
$\nu\ge s$, proving the first formula.  Under the no-crossing hypothesis,
$\chi_\nu=1$.  The full cubic sum is zero, while its three successive suffix
sums are $-1,2,-1$, proving the second formula.
\end{proof}

\begin{example}[cubic cancellation can fail with no floor crossings]
\label{long269:res:no-crossing-counterexample}
Take $a=0$, $r=35$ and $(c_0,c_1,c_2,c_3)=(1,-3,3,-1)$.  Exact integer
comparisons give
\[
\begin{array}{c|rrrr}
 u&0&35&70&105\\\hline
 \lfloor\log_3 2^u\rfloor=\lfloor\log_3 2^{u+1}\rfloor&0&22&44&66\\
 \lfloor\log_5 2^u\rfloor=\lfloor\log_5 2^{u+1}\rfloor&0&15&30&45\\
 m_u&1&195&723&1582
\end{array}
\]
The first two rows imply $b_u=2$ and
$\lfloor(u+t)\theta_p\rfloor=\lfloor u\theta_p\rfloor$ for every
$0\le t<1$, $p=3,5$.  Hence every relevant $\kappa_p(0,t,u)$ vanishes,
$\gamma_{0,u}=1$, and all triangle weights at these four indices are $1$.
The last row is therefore the count of pairs satisfying
$3^j5^k<2^{u+1}$, not a floating-point estimate.  The strips have sizes
$194,528,859$, so
\[
 D_{35,0}=15(-194+2\cdot528-859)=45\ne0.
\]
A portable certificate records the count in every $j$-row and verifies all
power inequalities using integers.  This single example disproves
unconditional cubic cancellation from carry avoidance.  It does not
exclude a different fixed operator, different shifts, a boundary correction
or a sparse-defect statement restricted to a suitable infinite subsequence.
\end{example}


\paragraph{What the Hecke--Mahler comparison does and does not supply.}
In the published Luca--Ouaknine--Worrell article
\cite{lucaouaknineworrell2025}, Definition~5 requires both expanding gaps
and polynomial variation of nonzero finite-difference defects.  Theorem~6
is the associated fixed-base value criterion; Theorem~8 and Claim~10
establish the hypotheses for polynomial floor sequences.  Claim~10 uses
exact finite-difference cancellation when no floor crossing occurs.
Our $m_a$ is instead a weighted lattice count.  Proposition~\ref{long269:res:strip-decomposition}
separates floor crossings from its moving boundary, and
Example~\ref{long269:res:no-crossing-counterexample} shows that the latter
can survive a cubic difference even with no crossings.  The correction
preserves the value, but does not prove sparseness.

A possible coefficient target is a fixed operator and shifts $r_n\to\infty$
for which $\Delta_n=\{a\ge0:D_{r_n,a}\ne0\}$ is infinite and distinct
members are at least $\eta r_n$ apart, for a fixed $\eta>0$.
Even this would leave the uniform polynomial-variation requirement, for
example
\[
 |D_{r_n,a'}|\le C\bigl((a'-a)^d+|D_{r_n,a}|\bigr)
 \quad(a<a',\ a,a'\in\Delta_n),
\]
with $C,d$ independent of $n$.
The bound $O((a+r+1)^2)$ is not this relative-variation estimate.
There is also a representation issue: the displayed denominator chain is
not a fixed-base power sequence.  The next proposition shows exactly what
is lost in two direct fixed-base recodings.  A theorem for the original
chain would need a separate proof, including non-cancellation and the
height estimates in any Subspace-Theorem argument.  The finite set of
prime divisors alone does not provide those estimates.

\paragraph{The echoing comparison has a different interface.}
Kebis, Luca, Ouaknine, Scoones and Worrell
\cite[Definition~3, Theorem~6 and Claim~7]{kebis2024echoing}
work with fixed-base series whose coefficients lie in a finite algebraic
alphabet, and require non-vanishing weighted mismatch sums.
The four-letter radix sequence is not the numerator sequence: the actual
$m_a$ is unbounded.  A recoding would have to preserve the scalar value,
identify its coefficient alphabet and verify those mismatch conditions.
No such recoding is proved here.

\begin{proposition}[what direct fixed-base recoding preserves]
\label{long269:res:fixed-base-recoding}
The following identities converge absolutely:
\[
 \alpha=\sum_{a\ge0}\frac{e_a}{30^{a+1}}
        =\sum_{a\ge0}\frac{v_a}{8^{a+1}},\qquad
 e_a=\frac{m_a30^{a+1}}{P_{a+1}},\quad
 v_a=\frac{m_a8^{a+1}}{P_{a+1}}.
\]
Here $e_a\in\mathbb Z_{>0}$ and $e_a\ge(15/4)^{a+1}$, whereas
$v_a\in\mathbb Z[1/15]$ and $0<v_a<225(a+1)^2$.
For an integer $q\ge2$, the termwise divisibility $P_n\mid q^n$ for every
$n\ge1$ holds exactly when $30\mid q$; that direct recoding then has
coefficients at least $(q/8)^{a+1}$.
\end{proposition}
\begin{proof}
Each exponent in $P_n$ is at most $n$, so $P_n\mid30^n$.  Also
$8^n/15<P_n\le8^n$ by the exact two-exponent and the strict three- and
five-power bounds.  These facts, $m_a\ge1$ and
$m_a\le15(a+1)^2$, give all the coefficient bounds and identities.
Necessity of $30\mid q$ follows because each of $2,3,5$ divides some $P_n$;
sufficiency follows from $P_n\mid30^n$.  The general lower bound follows
again from $P_n\le8^n$.
\end{proof}
The base-$30$ coefficients are integral but not polynomially bounded, so
they do not satisfy the hypotheses of the cited fixed-base polynomial-growth
criterion.  The base-$8$ coefficients have polynomial size, but integrality
and the required arithmetic height bounds are not established.  No claim
is made that their reduced denominators are unbounded: cancellation with
$m_a$ must be considered.  These are obstructions to the direct recodings,
not a proof that all possible fixed-base representations fail.
Adamczewski--Bugeaud's complexity theorem
\cite[Theorem~1]{adamczewskibugeaud2007} concerns the actual digits of an
integer-base expansion: for an algebraic irrational their length-$n$ block
complexity divided by $n$ tends to infinity.  Neither $(b_a)$ nor the
uncarried $(e_a)$ is such a digit expansion of $\alpha$.




\begin{problem}[the repeated three-prime target]
\label{long269:prob:producer}\label{long269:prob:tails269}
Prove irrationality of $S=\mathcal R_{\{2,3,5\}}$.  Equivalent forms are
$\Esc(K)$ and
\begin{equation}\label{long269:eq:tail-nonintegrality}
 BX_a\notin\mathbb Z
 \quad\text{for every }B\ge1\text{ with }\gcd(B,30)=1\text{ and every }a\ge1.
\end{equation}
\end{problem}
The first formulation uses finite integer data but unbounded denominator
and onset quantifiers; the second uses the actual real tail.  Neither is
an additional hypothesis imposed on the series, and neither is proved by
the rank obstruction.  The next proposition records their common scalar
target, complementing Theorem~\ref{long269:res:actual-escape-endpoint}.

\begin{proposition}[the reduced-tail question is also the target]
\label{long269:res:tails-equivalence}
Statement~\eqref{long269:eq:tail-nonintegrality}, quantified over every $B\ge1$
coprime to $30$ and every $a\ge1$, is equivalent to irrationality of $S$.
\end{proposition}

\begin{proof}
If $BX_a\in\Z$ for some such $B$ and $a$, then $X_a\in\Q$, and since
$S=\sum_{j<a}s_j+X_a/h_a$ with $h_a$ a positive rational and the prefix a
finite sum of rationals, $S\in\Q$.  Conversely if $S\in\Q$, then
Theorem~\ref{long269:res:actual-cancellation} produces $B$ coprime to $30$ with
$BX_a\in\Z$ for every $a\ge a_D$.
\end{proof}

Integrality at one index propagates forwards, but nonintegrality at one
early index does not propagate forever: a rational denominator can clear
later.  Proposition~\ref{long269:res:exact-denominator} describes that onset
exactly.  The scaled dichotomy also permits an integral branch, so it does
not settle the target.

\subsection{The analytic route and what it would need}

For the de-duplicated series put
\[
 \mathcal D_{2,3,5}=1+
 \sum_{t\in\{2^{n},3^{n},5^{n}:n\ge1\}}
 \frac{1}{2^{\lfloor\log_2t\rfloor}3^{\lfloor\log_3t\rfloor}
           5^{\lfloor\log_5t\rfloor}} ,
\]
whose dyadic coding is the joint rotation word
$\delta_{3,a}=\lfloor(a+1)\theta_3\rfloor-\lfloor a\theta_3\rfloor$ and
$\delta_{5,a}=\lfloor(a+1)\theta_5\rfloor-\lfloor a\theta_5\rfloor$, with
$b_a=2\cdot3^{\delta_{3,a}}5^{\delta_{5,a}}$.

\begin{problem}[function-faithful two-dimensional representation]
\label{long269:prob:representation}
The historical irrationality assertion for $\mathcal D_{2,3,5}$ is not
being reclassified as a new open problem.  A recovered proof or a
transcendence theorem would require its own argument.  For a functional
approach, first specify the function, its coefficient field and convergence
domain, then prove an exact identity for $\mathcal D_{2,3,5}$, and only then
apply a stated value theorem.  A conditional theorem must display the
extra nondegeneracy assumption; a no-representation theorem must specify
the class it excludes.

Individual irrationality of $\theta_3$ and $\theta_5$ is what
Theorem~\ref{long269:res:infinite-rank} uses, and it is weaker than the joint
equidistribution or the rational independence of $1,\theta_3,\theta_5$ that a
two-dimensional value theorem may need; the stronger hypothesis is not assumed
anywhere above.  The finite-observer formalisation isolates the precise
faithfulness requirement: equality in a finite observer must imply equality
after symbolic realisation, and a genuine finite-dimensional factorisation
forces the realised symbolic span to be finite-dimensional
(\lword{Erdos269/WeightedPhaseCarry.lean}{109}{carry_eq_residueDigit_add_coboundary}{residue-coboundary
form},
\lword{Erdos269/WeightedPhaseCarry.lean}{293}{CartierRealisation}{symbolic
realisation},
\lword{Erdos269/WeightedPhaseCarry.lean}{334}{finite_realisedSpan_of_factorisation}{checked}),
with the carry residue and residue digit confined to their declared intervals
(\lword{Erdos269/WeightedPhaseCarry.lean}{150}{carryResidue_mem_interval}{carry
interval},
\lword{Erdos269/WeightedPhaseCarry.lean}{157}{residueDigit_mem_interval}{digit
interval}).  No theorem here proves that the literal realised span is infinite.

\paragraph{Which value-theorem hypotheses remain missing.}
Pellarin's rank-one theorem
\cite[Theorem~1.1]{pellarin2006} treats Hecke--Mahler values with a
quadratic irrational slope and a specified rank-one module condition on
the algebraic evaluation points.  Two lattice coordinates alone imply
neither hypothesis.  For linear Mahler systems in one variable,
Adamczewski and Faverjon give a new proof of Nishioka's theorem
\cite[Theorems~1 and 2]{adamczewskifaverjon2023}; one must have a Mahler
functional system and a regular algebraic evaluation point.
Their multivariate theory is now published
\cite{adamczewskifaverjon2026}.  The author manuscript's Theorem~3.3 and
Sections~3 and~5 retain a genuine functional system, regularity of the
evaluation point and admissibility of the transformation--point pair.
No such system or evaluation data for the literal three-prime sum are
constructed in this record.  In particular, the theorem's freedom in the
system matrix does not remove its transformation and point hypotheses.
The source audit distinguishes the 68-page author manuscript used for
these locators, the older 52-page arXiv version and the journal metadata.

\end{problem}

\subsection{Structural constraints already in place}

The radix alphabet extends without difficulty.  For ordered primes
$p_1<\cdots<p_s$ an interval $(p_1^{a},p_1^{a+1})$ contains at most one power
from each other channel, because consecutive $p_i$-powers have ratio
$p_i>p_1$, so its block radix belongs to the $2^{s-1}$-letter alphabet
$\{p_1\prod_{i=2}^{s}p_i^{\varepsilon_i}:\varepsilon_i\in\{0,1\}\}$.  The
quantitative questions are the interesting ones: effective recurrence or
discrepancy for the actual four-letter $\{2,6,10,30\}$ word, an asymptotic with
an error term for the restricted two-dimensional shell counts that generate
$m_a$, or the exact separated rank of the literal kernel under a specified
family of shifts.

Three-channel rigidity and carry-lift extinction already exclude one proposed
argument in four exact steps.  Under channel surjectivity, ordinary block
nullity is equivalent to the perturbation being a coboundary of a channel
potential
(potential classifier).  Zero perturbations on genuine $2\to3$ and $2\to5$ transitions
then identify all three potential values and force the perturbation to vanish
at every index.  For an integral lift, that vanishing makes a nonzero initial
error grow by the exact product of the successive bases, and bases at least two
make its absolute value at least $2^{N}$, contradicting even a single
index-$N$ bound strictly below $2^{N}$
(one-index extinction); consequently no uniform bound on the lift error can hold either
(uniform extinction).  A separate four-state calculation reaches the obstruction
earlier: four real states in $(0,1)$ with unit-accuracy integral lifts and the
two anchor equalities force the first complete $2$-block sum to be $1$, so that
block cannot be null
(first-block sum,
four-state obstruction).

These lift conclusions are conditional.  The paper constructs the actual orbit
and its bounded integral carry under rationality, and it does not construct a
faithful lift with the two anchors or the block nullity those auxiliary
statements require.  The actual carry supplies a weighted block defect instead,
so any successful argument along that line must use the weighted identity or
construct a different faithful lift.

A single conditional single-channel criterion is also on record.  For the pure
$2$-channel sum the Cantor states lie in an explicit open interval
(confinement),
consecutive states are separated
(gap bound), a small nonzero linear form forces irrationality
(criterion),
and clearing together with small gaps assembles to irrationality
(assembly).
The clearing and small-gap hypotheses are not discharged for any concrete
series, so no irrationality statement follows from them here.

\subsection{Where the problem stands}

The two-prime transcendence deduction is an earlier application of the
cited value theorem.  For the repeated $\{2,3,5\}$ value, this record
establishes the literal coefficients, integral recurrence, exact clearing,
quadratic bounds and the equivalent residue target.  It does not establish
irrationality.  The rank and uniform-norm obstructions concern the kernel,
not the arithmetic of its sum.  The fixed-start residue formula clarifies
the target's quantifiers, while the strip decomposition identifies a
concrete obstruction to transplanting cubic no-crossing cancellation.
A further argument must use the actual weighted multiplicities to exclude
eventual integral reduced tails.  Finite tests, countability, low-complexity
radix coding and denominator support do not by themselves do so.

\pdfbookmark[1]{Statements and declarations}{statements-declarations}
\subsection*{Statements and declarations}

\paragraph{Proof sources.}
The two-prime deduction rests on the cited external value theorem and is
not formalised here; Fan's recorded priority is retained.
The attached \texttt{LEAN\_INDEX.json} at public commit
\nolinkurl{3d6d938d696fed0fb71dd55115a18a73738ff223} records a successful CI
build for the relevant public source closure.  This includes arbitrary-order
kernel rank, the actual tail recurrence, rationality clearing, the actual
$Q$ and $\widetilde Q$ bounds, and the $o(8^a)$ equivalence band.
The separate Palomar release at \texttt{52f29ad1} selects the
arbitrary-order uniform-minor and prime non-separation clauses, as well as
the finite fixture.  Its Comparator evidence is therefore not limited to a
$2\times2$ determinant.
These statements report the supplied receipts; this editorial revision did
not rerun Lean, Isabelle, Comparator or NanoDa.  The source index is not a
blanket axiom audit.  Individual historical links keep their own pins and
must not be treated as byte-identical to the newer tree without a comparison.
The weighted-triangle and weighted-shift identities, exact denominator
formula, scaled dichotomy, strengthened rigidity, residue limit, strip
decomposition and direct recoding proposition have ordinary proofs.
The new finite checks support their finite instances, not their infinite
quantifiers.  None is presented as a new formalisation result of this pass.
The large numerical exclusions remain archived reports until their witnesses
are published.  No evidence item supplies the source-specific cofinal escape.

\paragraph{Artefact and data availability.} The
\href{\sourceurl}{pinned formal-source revision} contains the Lean sources, the
fixed toolchain, the library manifest, and the exact dyadic-window checker used
in the finite scan.  Those sources support only the formal statements
identified in the evidence register.  Ordinary proofs and external analytic
inputs have their own stated dependencies; the unavailable numerical
witnesses are not supplied by a repository link.

\paragraph{Funding and competing interests.} This work received no external
funding.  The author declares no competing interests.

\paragraph{Acknowledgements.} Priority for the two-prime factorisation, its
Hecke--Mahler reduction and the running-LCM identity for every $|P|$ belongs to
Steve Fan's forum post of 26 June 2026~\cite{fan2026comment}.  The
transcendence input is due to Yann Bugeaud and Michel Laurent and to the
earlier work of Loxton and van der Poorten cited by them.  The
problem numbering and historical status snapshot are taken from the
Erd\H{o}s Problems catalogue maintained by Thomas Bloom~\cite{erdosproblems}.
